\documentclass{article}
\usepackage{ifluatex}
\ifluatex 
    \usepackage{fontspec}
    \setsansfont{CMU Sans Serif}%{Arial}
    \setmainfont{CMU Serif}%{Times New Roman}
    \setmonofont{CMU Typewriter Text}%{Consolas}
    \defaultfontfeatures{Ligatures={TeX}}
\else
    \usepackage[T2A]{fontenc}
    \usepackage[utf8]{inputenc}
\fi
\usepackage[english,russian]{babel}
\usepackage{amssymb,latexsym,amsmath,amscd,mathtools,wasysym}
\usepackage[shortlabels]{enumitem}
\usepackage[makeroom]{cancel}
\usepackage{graphicx}
\usepackage{geometry}
\usepackage{verbatim}
\usepackage{fvextra}
\usepackage{multicol}
\usepackage{multirow}
\usepackage{tabu}
\usepackage{arydshln} % \hdashline and :
\usepackage{float}
\makeatletter
\g@addto@macro\@floatboxreset\centering
\makeatother
\usepackage{caption}
\usepackage{csquotes}
\usepackage[bb=dsserif]{mathalpha}
\usepackage[normalem]{ulem}

\usepackage{xcolor}
\colorlet{darkgreen}{black!25!blue!50!green}


%% Here f*cking with mathabx
\DeclareFontFamily{U}{matha}{\hyphenchar\font45}
\DeclareFontShape{U}{matha}{m}{n}{
    <5> <6> <7> <8> <9> <10> gen * matha
    <10.95> matha10 <12> <14.4> <17.28> <20.74> <24.88> matha12
}{}
\DeclareSymbolFont{matha}{U}{matha}{m}{n}
\DeclareFontFamily{U}{mathb}{\hyphenchar\font45}
\DeclareFontShape{U}{mathb}{m}{n}{
    <5> <6> <7> <8> <9> <10> gen * mathb
    <10.95> matha10 <12> <14.4> <17.28> <20.74> <24.88> mathb12
}{}
\DeclareSymbolFont{mathb}{U}{mathb}{m}{n}

\DeclareMathSymbol{\defeq}{\mathrel}{mathb}{"15}
\DeclareMathSymbol{\eqdef}{\mathrel}{mathb}{"16}


\usepackage{trimclip}
\DeclareMathOperator{\updownarrows}{\clipbox{0pt 0pt 4.175pt 0pt}{$\upuparrows$}\hspace{-.825px}\clipbox{0pt 0pt 4.175pt 0pt}{$\downdownarrows$}}
\DeclareMathOperator{\downuparrows}{\clipbox{0pt 0pt 4.175pt 0pt}{$\downdownarrows$}\hspace{-.825px}\clipbox{0pt 0pt 4.175pt 0pt}{$\upuparrows$}}

\makeatletter
\providecommand*\deletecounter[1]{%
    \expandafter\let\csname c@#1\endcsname\@undefined}
\makeatother

\usepackage{hyperref}
\hypersetup{
    colorlinks=true,
    linkcolor=cyan,
    urlcolor=blue,
    breaklinks=true,
}

\usepackage{tikz,tikz-3dplot}
\usepackage{pgfplots}
\pgfplotsset{compat=newest}
\usepgfplotslibrary{fillbetween,polar}

\usepackage{tkz-euclide}
\usetikzlibrary{calc,patterns,angles,quotes,backgrounds,shapes.geometric,trees,positioning}

\ifluatex
\usetikzlibrary{graphs,graphs.standard,graphdrawing}
\usegdlibrary{trees,force}
\else
\errmessage{Run with LuaTeX, if you want to use gdlibraries}
\fi
\newcommand{\convexpath}[2]{
    [   
    create hullcoords/.code={
        \global\edef\namelist{#1}
        \foreach [count=\counter] \nodename in \namelist {
            \global\edef\numberofnodes{\counter}
            \coordinate (hullcoord\counter) at (\nodename);
        }
        \coordinate (hullcoord0) at (hullcoord\numberofnodes);
        \pgfmathtruncatemacro\lastnumber{\numberofnodes+1}
        \coordinate (hullcoord\lastnumber) at (hullcoord1);
    },
    create hullcoords
    ]
    ($(hullcoord1)!#2!-90:(hullcoord0)$)
    \foreach [
    evaluate=\currentnode as \previousnode using \currentnode-1,
    evaluate=\currentnode as \nextnode using \currentnode+1
    ] \currentnode in {1,...,\numberofnodes} {
        let \p1 = ($(hullcoord\currentnode) - (hullcoord\previousnode)$),
        \n1 = {atan2(\y1,\x1) + 90},
        \p2 = ($(hullcoord\nextnode) - (hullcoord\currentnode)$),
        \n2 = {atan2(\y2,\x2) + 90},
        \n{delta} = {Mod(\n2-\n1,360) - 360}
        in 
        {arc [start angle=\n1, delta angle=\n{delta}, radius=#2]}
        -- ($(hullcoord\nextnode)!#2!-90:(hullcoord\currentnode)$) 
    }
} % Nodes CW. Supposes tikz >= 3.0, else swap atan2 arguments.

\makeatletter
\newcommand*\circled[2][1.0]{\tikz[baseline=(char.base)]{
        \node[shape=circle, draw, inner sep=2pt,
        minimum height={\f@size*#1},] (char) {#2};}}
\makeatother

\newcommand{\existence}{{\circled{$\exists$}}}
\newcommand{\uniqueness}{{\circled{$\hspace{0.5px}!$}}}
\newcommand{\rightimp}{{\circled{$\Rightarrow$}}}
\newcommand{\leftimp}{{\circled{$\Leftarrow$}}}

\DeclareMathOperator{\sign}{sign}
\DeclareMathOperator{\Cl}{Cl}
\DeclareMathOperator{\proj}{pr}
\DeclareMathOperator{\Arg}{Arg}
\DeclareMathOperator{\supp}{supp}
\DeclareMathOperator{\diag}{diag}
\DeclareMathOperator{\tr}{tr}
\DeclareMathOperator{\rank}{rank}
\DeclareMathOperator{\Lat}{Lat}
\DeclareMathOperator{\Lin}{Lin}
\DeclareMathOperator{\Ln}{Ln}
\DeclareMathOperator{\Orbit}{Orbit}
\DeclareMathOperator{\St}{St}
\DeclareMathOperator{\Seq}{Seq}
\DeclareMathOperator{\PSet}{PSet}
\DeclareMathOperator{\MSet}{MSet}
\DeclareMathOperator{\Cyc}{Cyc}
\DeclareMathOperator{\Hom}{Hom}
\DeclareMathOperator{\End}{End}
\DeclareMathOperator{\Aut}{Aut}
\DeclareMathOperator{\Ker}{Ker}
\DeclareMathOperator{\Def}{def}
\DeclareMathOperator{\Alt}{Alt}
\DeclareMathOperator{\Sim}{Sim}
\DeclareMathOperator{\Int}{Int}
\DeclareMathOperator{\grad}{grad}
\DeclareMathOperator{\sech}{sech}
\DeclareMathOperator{\csch}{csch}
\DeclareMathOperator{\asin}{\sin^{-1}}
\DeclareMathOperator{\acos}{\cos^{-1}}
\DeclareMathOperator{\atan}{\tan^{-1}}
\DeclareMathOperator{\acot}{\cot^{-1}}
\DeclareMathOperator{\asec}{\sec^{-1}}
\DeclareMathOperator{\acsc}{\csc^{-1}}
\DeclareMathOperator{\asinh}{\sinh^{-1}}
\DeclareMathOperator{\acosh}{\cosh^{-1}}
\DeclareMathOperator{\atanh}{\tanh^{-1}}
\DeclareMathOperator{\acoth}{\coth^{-1}}
\DeclareMathOperator{\asech}{\sech^{-1}}
\DeclareMathOperator{\acsch}{\csch^{-1}}

\newcommand*{\scriptA}{{\mathcal{A}}}
\newcommand*{\scriptB}{{\mathcal{B}}}
\newcommand*{\scriptC}{{\mathcal{C}}}
\newcommand*{\scriptD}{{\mathcal{D}}}
\newcommand*{\scriptE}{{\mathcal{E}}}
\newcommand*{\scriptF}{{\mathcal{F}}}
\newcommand*{\scriptH}{{\mathcal{H}}}
\newcommand*{\scriptK}{{\mathcal{K}}}
\newcommand*{\scriptL}{{\mathcal{L}}}
\newcommand*{\scriptM}{{\mathcal{M}}}
\newcommand*{\scriptP}{{\mathcal{P}}}
\newcommand*{\scriptQ}{{\mathcal{Q}}}
\newcommand*{\scriptR}{{\mathcal{R}}}
\newcommand*{\scriptT}{{\mathcal{T}}}
\newcommand*{\scriptU}{{\mathcal{U}}}
\newcommand*{\scriptX}{{\mathcal{X}}}
\newcommand*{\Cnk}[2]{\left(\begin{matrix}#1\\#2\end{matrix}\right)}
\newcommand*{\im}{{\mathbf i}}
\newcommand*{\id}{{\mathrm{id}}}
\newcommand*{\compl}{^\complement}
\newcommand*{\trans}{^{\mathsf T}}
\newcommand*{\dotprod}[2]{{\left\langle{#1},{#2}\right\rangle}}
\newcommand\matr[1]{\left(\begin{matrix}#1\end{matrix}\right)}
\newcommand\matrd[1]{\left|\begin{matrix}#1\end{matrix}\right|}
\newcommand\arr[2]{\left(\begin{array}{#1}#2\end{array}\right)}
\newcommand{\dfn}{\item[\color{red}\textbullet]}
\newcommand{\thm}{\item[\color{darkgreen}\textbullet]}

\DeclareMathOperator{\divby}{\scalebox{1}[.65]{\vdots}}

\newcommand{\undercolorblack}[2]{{\color{#1}\underline{\color{black}#2}}}
\newcommand{\undercolor}[2]{{\colorlet{tmp}{.}\color{#1}\underline{\color{tmp}#2}}}

\usepackage{adjustbox}

\usepackage{minted}
\usemintedstyle{tango}

\tikzset{treenode/.style={
    inner sep=0pt,
    text width=5mm,
    scale=.8,
    circle,
    outer sep=2mm,
    align=center,
    draw=black,
    fill=white,
    thin
}}

\geometry{legalpaper, paperheight=16383pt, margin=1in}
\setcounter{totalnumber}{100}
\pagestyle{empty}

\newcommand{\pseudo}{PseudoLexer.py:PseudoLexer -x}
\newcommand{\pseudocode}[1]{\mintinline{\pseudo}{#1}}

\begin{document}
    \paragraph{\undercolorblack{orange}{Алгоритмы, оценка сложности, прочая штука}}
    Итак, что такое алгоритм? Ну, фиг его знает так-то. Но вообще это какая-то, обычно конечная, последовательность действий, которая выполняет какую-то формальную задачу.
    Обычно, у него есть какой-то вход и какой-то выход. По жизни может быть хуже.\\
    Пример задачи: задан массив, надо посчитать его минимум. Псевдокод. Псевдокод --- это кто? Псевдокод --- это друг человека.
    \begin{minted}{\pseudo}
        ans = inf
        for i in [0, n-1]
            ans = min(ans, a[i])
        return ans
    \end{minted}
    Теперь хочется как-то понять, а этот алгоритм, он \textit{хороший} --- или нет?\\
    Временная сложность алгоритма --- это количество действий, который он совершает. Осталось лишь понять, что такое \textit{действие}. Обычно мы задаём, что у нас некоторые штуки, такие как \pseudocode{+}, \pseudocode{-} и ещё какие-нибудь --- это одно действие, а всё остальное выражается через них. Такой набор действий называется моделью вычислений. Самая простая из них --- RAM-модель. Согласно ей, все данные лежат в одном массиве, и любое обращение к этому массиву --- одно действие. Вообще говоря, для обращения к жёстким дискам нужна другая модель, если у нас развита параллельность, --- третья и т.д. Но нам это слабо интересно. Также в RAM-модели есть сложные операции вида \pseudocode{for} или ещё чего, но мы формально это задавать не будем, предполагая, что мы знаем, что это и как оно работает.\\
    Вернёмся к нашему алгоритму. Сколько операций он совершает?
    \begin{minted}{\pseudo}
        ans = inf                      # Ну, одно.
        for i in [0, n-1]          # Ну, 2*n.
            ans = min(ans, a[i]) # Ну, 3*n.
        return ans                    # Ну, 1.
    \end{minted}
    Итого \pseudocode{5*n+2}. Довольно примерная оценка, но нам и не нужны точные, мы рассматриваем \textit{асимптотические оценки}. Говорят, что $f(n)=O(g(n))$, если $\exists N_0,c>0~\forall N\geqslant N_0~f(n)\leqslant cg(n)$. Физический смысл этой штуки: <<Начиная с какого-то момента $f(n)$ отличается от $g(n)$ не более, чем в $c$ раз.>> Про $O(n)$ есть несколько тупых утверждений. Например, такое: $a_kn^k+\cdots+a_in^i+\cdots+a_1n^1+a_0n^0=O(n^k)$.\\
    Говорят, что $f(n)=\varOmega(g(n))$, если $\exists N_0,c>0~\forall N\geqslant N_0~f(n)\geqslant cg(n)$. Если $f(n)=O(g(n))$ и $f(n)=\varOmega(g(n))$, то пишут, что $f(n)=\varTheta(g(n))$. И культурные люди пишут про оценку сложности алгоритмов $\varTheta$. Но мы не культурные, и будем писать $O$ по следующей причине: нам обычно не нужно оценивать алгоритмы снизу. Фраза <<мой алгоритм работает быстрее, чем>> полезнее, чем <<мой алгоритм работает медленнее, чем>>.
    \subparagraph{\undercolorblack{orange}{Сортировки}.}
    Если у нас есть отношение порядка, то можно сортировать.\\
    Сортировка вставками работает так. Мы ходим по массиву от \pseudocode{0} до \pseudocode{n-1} и на каждом этапе хотим поддерживать следующий инвариант: подмассив \pseudocode{[0:i]} отсортированным. Для этого мы каждый новый найденный элемент вставляем туда, куда его надо вставить. То есть этот элемент \pseudocode{a[i]} мы двигаем влево, пока это корректно. Как выглядит код:
    \begin{minted}{\pseudo}
        for i in [0..n-1]
            j = i
            while j > 0 && a[j] < a[j-1]
                swap(a[j], a[j-1])
                j--
    \end{minted}
    Какова сложность этой штуки? В самом плохом случае (когда массив отсортирован в обратном порядке) --- $O(n^2)$. В самом хорошем (когда отсортирован как надо) --- $O(n)$.\\
    И тут возникает мем. Мы почти всегда оцениваем худшее время по двум причинам. Во-первых, это гарантирует, что алгоритм не может работать уж совсем плохо. А во-вторых мы слишком слабы, чтобы в общем случае оценить, за какое время работает алгоритм, то есть выяснить функцию $T(a_1,a_2,\ldots,a_n)$. Тут мы ещё можем посчитать --- это $O(n+\text{количество инверсий в массиве \pseudocode{a}})$, а вот потом пойдёт полный капец.\\
    Что ж, как мы оцениваем время? С \pseudocode{for}'ами всё понятно. С \pseudocode{while}'ами уже потруднее, там нужно думать, но ещё терпимо. А вот с рекурсией нужно мучиться. Например:
    \begin{minted}{\pseudo}
        f(n):
            if n == 1
                return 1
            else
                return f(n-1)
    \end{minted}
    Такая штука работает за $O(n)$.
    \begin{minted}{\pseudo}
        f(n):
            if n == 1
                return 1
            else
                return f(n / 2)
    \end{minted}
    Такая --- за $O(\log n)$.
    \begin{minted}{\pseudo}
        f(n):
            if n == 1
                return 1
            else
                return f(n / 2) + f(n / 2)
    \end{minted}
    Чтобы понять, за сколько работает это, нало осознать, что дерево рекурсии тут двоичное и имеет высоту $\log_2n$, и на каждом этапе вдвое больше вызовов, чем на предыдущем. Значит всего эта программа выполняется за $2^{\log_2(n)}=n$. Совершенно аналогично, такая программа:
    \begin{minted}{\pseudo}
        f(n):
            if n == 1
                return 1
            else
                return f(n / 2) + f(n / 2) + f(n / 2)
    \end{minted}
    Работает за $3^{\log_2n}=n^{\log_23}$.\\
    К чему всё это? А вот к сортировке слиянием. Идея её состоит в том, что мы делам наш массив на две половинки, рекурсивно сортируем каждую, после чего соединяем. Как соединяем? А вот так:
    \begin{minted}{\pseudo}
        merge(a, b):
            i = j = 0
            while i < len(a) || j < len(b)
                if j == len(b) || i < len(a) && a[i] < b[j]
                    c[i+j] = a[i++]
                else
                    c[i+j] = b[j++]
            return c

        sort(c):
            if len(c) <= 1
                return c
            l = c[:len(c)/2-1]
            r = c[len(c)/2:]
            l = sort(l)
            r = sort(r)
            return merge(l, r)
    \end{minted}
    За какое время работает \pseudocode{merge}? Сумму \pseudocode{len(a)} и \pseudocode{len(b)}. То есть, когда мы вызываем её в \pseudocode{sort}, за $O(n)$. Все нерекурсивные операции работают за $O(1)$. осталось понять только что там с рекурсией. Если функцию количество операций данного алгоритма обозначит за $T$, то получим следующее равенство $T(n)=2T(\frac n2)+n$. Получается, что $T(n)=O(n\log n)$. Можно доказать это по индукции по честному, но мы не будем.
    \subparagraph{\undercolorblack{orange}{Мастер-теорема}.}
    Мастер-теорема заключается в следующем: Если $T(n)=aT\left(\frac nb\right)+n^k$, то
    \begin{itemize}
        \item При $k>\log_ba$ $T(n)=O(n^k)$.
        \item При $k=\log_ba$ $T(n)=O(n^k\log n)$.
        \item При $k<\log_ba$ $T(n)=O(n^{\log_ba})$.
    \end{itemize}
    Пример пользования этой штукой. Пусть нам надо что-то умножить, например, два длинных числа. Или два многочлена. Или ещё чего. Тогда алгоритм Карацубы выглядит так. Для примера мы будем умножать два 32-битных числа. Мы бьём наши множители (назовём их $a$ и $b$) на две половинки: $a_0$, $a_1$, $b_0$, $b_1$. Соответственно, $a=2^{16}\cdot a_0+a_1$, $b=2^{16}\cdot b_0+b_1$. Тогда $a\cdot b=2^{32}\cdot a_0b_0+2^{16}\cdot(a_0b_1+a_1b_0)+a_1b_1$. Если честно рекурсивно посчитать произведение четырёх пар 16-битных чисел, то получится $T(n)=4T\left(\frac n2\right)+n$, то есть $O(n^2)$. Это и руками получить можно. Но есть хак. Мы можем посчитать $(a_0+a_1)(b_0+b_1)$. Сложение, надо сказать, работает быстро, на него мы в асимптотике забиваем. При этом, поскольку $a_0b_0$ и $a_1b_1$ мы всё равно считаем, можно взять $(a_0+a_1)(b_0+b_1)-a_0b_0-a_1b_1$, что и будет равно коэффициенту при $2^{16}$. А значит тут уже получается, что $T(n)=3T\left(\frac n2\right)+n$, то есть $O(n^{\log_23})$. Примерно по такой же механике работает алгоритм Штрассена для перемножения матриц.
    \paragraph{\undercolorblack{orange}{Структуры данных}.}
    Это структура, хранящая данные. Первый вопрос --- зачем? Чтобы быстрее работали какие-то операции, которые нам хочется делать. Поэтому, если операции у нас разные, структуры данных тоже могут быть разными.\\
    Вот есть у нас массив. В нём быстро выполняются операции взятия и задания элемента: \pseudocode{get(i)} и \pseudocode{set(i, val)}. Обе операции работают за сколько? За $O(1)$.\\
    Почему СД и алгоритмы --- один курс? Потому что эти штуки очень связаны, в сложных алгоритмах может использоваться структуры данных и наоборот.\\
    \subparagraph{\undercolorblack{orange}{Куча}.}
    Хранит мультимножество линейно упорядоченных элементов. На ней есть две с половиной операции: \pseudocode{insert(x)} --- добавляет элемент в кучу, \pseudocode{get_min()} --- возвращает минимальный элемент кучи, \pseudocode{remove_min()} --- удаляет один из минимальных элементов кучи.\\
    Как это реализовать по-тупому? Массивом. У нас есть массив бесконечного размера. В левом его куске хранятся данные с индексами от \pseudocode{0} до \pseudocode{n-1}. Как добавить в него элемент:
    \begin{minted}{\pseudo}
        insert(x):
            a[n++] = x
    \end{minted}
    Как искать минимум?
    \begin{minted}{\pseudo}
        get_min():
            res = inf
            for i in [0:n-1]
                res = min(res, a[i])
            return res
    \end{minted}
    Как удалять минимальный элемент?
    \begin{minted}{\pseudo}
        get_min():
            j = 0
            for i in [0:n-1]
                if a[i] < a[j]:
                    j = i
            swap(a[j], a[--n])
    \end{minted}
    Измерим время работы всего. \pseudocode{insert} --- $O(1)$, \pseudocode{get_min} и \pseudocode{remove_min} --- $O(n)$.\\
    Идея номер два --- хранить индекс минимума. Назовём его \pseudocode{m}. Тогда
    \begin{minted}{\pseudo}
        insert(x):
            if a[m] > x
                m = n
            a[n++] = x

        get_min():
            return a[m]

        get_min():
            swap(a[m], a[--n])
            m = 0
            for i in [0:n-1]
                if a[i] < a[m]:
                    m = i
    \end{minted}
    Теперь \pseudocode{get_min} работает за $O(1)$.\\
    Идея номер 3. Хранить отсортированный по убыванию массив.
    \begin{minted}{\pseudo}
        remove_min():
            n--

        get_min():
            return a[n-1]

        insert(x):
            a[n++] = x
            j = n - 1
            while j > 0 && a[j] > a[j-1]
                swap(a[j], a[j-1])
                j--
    \end{minted}
    Теперь у нас очень быстро ($O(1)$) работают \pseudocode{get_min} и \pseudocode{remove_min}, но медленно ($O(n)$) работает \pseudocode{insert}.\\
    Вопрос: какая версия лучше, первая или вторая? См. аксиому Эскобара. Вообще можно скрещивать структуры. Конкретно эти две скрещиваются так: возьмём много маленьких структур данных (размера $b$), в каждой из которых хранится отсортированный массив. При добавлении элемента мы кладём элемент в последнюю маленькую кучу. Туда, куда надо. Это выполняется за $O(b)$. Чтобы найти и/или удалить минимум можно просмотреть минимумы маленьких структур, получится $O(\frac nb)$. Если мы хотим, чтобы эти операции выполнялись за одинаковое время, можно взять $b\approx\sqrt n$, тогда получится.\\
    Перейдём к правильной человеческой \textit{двоичной куче}. Построим двоичное дерево. Причём все слои, кроме последнего, заполнены. Наше дерево будет хранится не так, кок нормально, потому что полное дерево на какое-то количество вершин всегда одинаковое. Например, дерево на 12 вершин всегда выглядит так:\\
    \begin{tikzpicture}[
        scale=.5
    ]
        \node[treenode] (root) at (0,0) {};
        \node[treenode] (l) at (-4,-2) {};
        \node[treenode] (r) at (4,-2) {};
        \node[treenode] (ll) at (-6,-4) {};
        \node[treenode] (lr) at (-2,-4) {};
        \node[treenode] (rl) at (2,-4) {};
        \node[treenode] (rr) at (6,-4) {};
        \node[treenode] (lll) at (-7,-6) {};
        \node[treenode] (llr) at (-5,-6) {};
        \node[treenode] (lrl) at (-3,-6) {};
        \node[treenode] (lrr) at (-1,-6) {};
        \node[treenode] (rll) at (1,-6) {};
        
        \draw[->] (root) -- (l);
        \draw[->] (root) -- (r);
        \draw[->] (l) -- (ll);
        \draw[->] (l) -- (lr);
        \draw[->] (r) -- (rl);
        \draw[->] (r) -- (rr);
        \draw[->] (ll) -- (lll);
        \draw[->] (ll) -- (llr);
        \draw[->] (lr) -- (lrl);
        \draw[->] (lr) -- (lrr);
        \draw[->] (rl) -- (rll);
    \end{tikzpicture}\\
    % Сюда вставить картинку дерева на 12 элементов с элементами. Можно ещё массив нарисовать.
    Пронумеруем элементы так: корень --- \pseudocode{0}, его дети \pseudocode{1} и \pseudocode{2}, дети \pseudocode{1} --- \pseudocode{3} и \pseudocode{4} и так далее. Потом положим в массив по заданным индексам заданные элементы. Заметим, что левый ребёнок вершины с номером \pseudocode{i} имеет номер \pseudocode{2 * i + 1}, а правый --- \pseudocode{2 * i + 2}. А его родитель --- \pseudocode{(i - 1) / 2}. Дальше разложим элементы в дереве так: в любом поддереве его минимум лежит в корне:\\
    \begin{tikzpicture}[
        scale=.5
    ]
        \node[treenode] (root) at (0,0) {2};
        \node[treenode] (l) at (-4,-2) {5};
        \node[treenode] (r) at (4,-2) {7};
        \node[treenode] (ll) at (-6,-4) {8};
        \node[treenode] (lr) at (-2,-4) {6};
        \node[treenode] (rl) at (2,-4) {10};
        \node[treenode] (rr) at (6,-4) {42};
        \node[treenode] (lll) at (-7,-6) {11};
        \node[treenode] (llr) at (-5,-6) {15};
        \node[treenode] (lrl) at (-3,-6) {28};
        \node[treenode] (lrr) at (-1,-6) {9};
        \node[treenode] (rll) at (1,-6) {13};
        
        \draw[->] (root) -- (l);
        \draw[->] (root) -- (r);
        \draw[->] (l) -- (ll);
        \draw[->] (l) -- (lr);
        \draw[->] (r) -- (rl);
        \draw[->] (r) -- (rr);
        \draw[->] (ll) -- (lll);
        \draw[->] (ll) -- (llr);
        \draw[->] (lr) -- (lrl);
        \draw[->] (lr) -- (lrr);
        \draw[->] (rl) -- (rll);
    \end{tikzpicture}\hspace{1.5cm}
    \begin{tikzpicture}[scale=.65, every node/.style={scale=0.67}]
        \draw[step=1cm,black,very thin] (0,0) grid (12,1);
        \node[scale=2] at (.5,.5) {2};
        \node[scale=2] at (1.5,.5) {5};
        \node[scale=2] at (2.5,.5) {7};
        \node[scale=2] at (3.5,.5) {8};
        \node[scale=2] at (4.5,.5) {6};
        \node[scale=2] at (5.5,.5) {10};
        \node[scale=2] at (6.5,.5) {42};
        \node[scale=2] at (7.5,.5) {11};
        \node[scale=2] at (8.5,.5) {15};
        \node[scale=2] at (9.5,.5) {28};
        \node[scale=2] at (10.5,.5) {9};
        \node[scale=2] at (11.5,.5) {13};
    \end{tikzpicture}\\
    Как в таком случае выглядят операции?
    \begin{minted}{\pseudo}        
        get_min():
            return a[0]

        insert(x):
            # Inserting to the first free space.
            a[n++] = x
            # Sifting up.
            i = n - 1
            while i != 0 && a[i] < a[(i - 1) / 2]
                swap(a[i], a[(i - 1) / 2])
                i = (i - 1) / 2

         remove_min():
            # Moving last element to the root.
            a[0] = a[--n]
            # Sifting down
            i = 0
            while 2 * i + 1 < n
                # j is minimal existing child
                j = 2 * i + 1
                if 2 * i + 2 < n && a[2 * i + 2] < a[2 * i + 1]
                    j = 2 * i + 2
                if a[i] <= a[j]
                    break
                swap(a[i], a[j])
                i = j
    \end{minted}
    Какая сложность \pseudocode{insert}? $O(\log n)$. Потому что у нас дерево ровное, красивое, мы совершаем не более, чем по 2 операции, на каждом слое (которых $\log_2n$). Аналогично \pseudocode{remove_min}.\\
    Иллюстрация \pseudocode{insert}:\\
    \begin{tikzpicture}[
        scale=.4,
        every node/.style={scale=.8}
    ]
        \node[treenode] (root) at (0,0) {2};
        \node[treenode] (l) at (-4,-2) {5};
        \node[treenode] (r) at (4,-2) {7};
        \node[treenode] (ll) at (-6,-4) {8};
        \node[treenode] (lr) at (-2,-4) {6};
        \node[treenode] (rl) at (2,-4) {10};
        \node[treenode] (rr) at (6,-4) {42};
        \node[treenode] (lll) at (-7,-6) {11};
        \node[treenode] (llr) at (-5,-6) {15};
        \node[treenode] (lrl) at (-3,-6) {28};
        \node[treenode] (lrr) at (-1,-6) {9};
        \node[treenode] (rll) at (1,-6) {13};
        \node[treenode, rectangle, minimum height=5mm, text centered] (rlr) at (3,-6) {4};
        
        \draw[->] (root) -- (l);
        \draw[->] (root) -- (r);
        \draw[->] (l) -- (ll);
        \draw[->] (l) -- (lr);
        \draw[->] (r) -- (rl);
        \draw[->] (r) -- (rr);
        \draw[->] (ll) -- (lll);
        \draw[->] (ll) -- (llr);
        \draw[->] (lr) -- (lrl);
        \draw[->] (lr) -- (lrr);
        \draw[->] (rl) -- (rll);
        \draw[->, red] (rl) -- (rlr);
    \end{tikzpicture}
    \begin{tikzpicture}[
        scale=.4,
        every node/.style={scale=.8}
    ]
        \node[treenode] (root) at (0,0) {2};
        \node[treenode] (l) at (-4,-2) {5};
        \node[treenode] (r) at (4,-2) {7};
        \node[treenode] (ll) at (-6,-4) {8};
        \node[treenode] (lr) at (-2,-4) {6};
        \node[treenode, rectangle, minimum height=5mm, text centered] (rl) at (2,-4) {4};
        \node[treenode] (rr) at (6,-4) {42};
        \node[treenode] (lll) at (-7,-6) {11};
        \node[treenode] (llr) at (-5,-6) {15};
        \node[treenode] (lrl) at (-3,-6) {28};
        \node[treenode] (lrr) at (-1,-6) {9};
        \node[treenode] (rll) at (1,-6) {13};
        \node[treenode] (rlr) at (3,-6) {10};
        
        \draw[->] (root) -- (l);
        \draw[->] (root) -- (r);
        \draw[->] (l) -- (ll);
        \draw[->] (l) -- (lr);
        \draw[->, red] (r) -- (rl);
        \draw[->] (r) -- (rr);
        \draw[->] (ll) -- (lll);
        \draw[->] (ll) -- (llr);
        \draw[->] (lr) -- (lrl);
        \draw[->] (lr) -- (lrr);
        \draw[->] (rl) -- (rll);
        \draw[->] (rl) -- (rlr);
    \end{tikzpicture}
    \begin{tikzpicture}[
        scale=.4,
        every node/.style={scale=.8}
    ]
        \node[treenode] (root) at (0,0) {2};
        \node[treenode] (l) at (-4,-2) {5};
        \node[treenode, rectangle, minimum height=5mm, text centered] (r) at (4,-2) {4};
        \node[treenode] (ll) at (-6,-4) {8};
        \node[treenode] (lr) at (-2,-4) {6};
        \node[treenode] (rl) at (2,-4) {7};
        \node[treenode] (rr) at (6,-4) {42};
        \node[treenode] (lll) at (-7,-6) {11};
        \node[treenode] (llr) at (-5,-6) {15};
        \node[treenode] (lrl) at (-3,-6) {28};
        \node[treenode] (lrr) at (-1,-6) {9};
        \node[treenode] (rll) at (1,-6) {13};
        \node[treenode] (rlr) at (3,-6) {10};
        
        \draw[->] (root) -- (l);
        \draw[->, darkgreen] (root) -- (r);
        \draw[->] (l) -- (ll);
        \draw[->] (l) -- (lr);
        \draw[->] (r) -- (rl);
        \draw[->] (r) -- (rr);
        \draw[->] (ll) -- (lll);
        \draw[->] (ll) -- (llr);
        \draw[->] (lr) -- (lrl);
        \draw[->] (lr) -- (lrr);
        \draw[->] (rl) -- (rll);
        \draw[->] (rl) -- (rlr);
    \end{tikzpicture}\\
    Иллюстрация \pseudocode{remove_min} (напоминаю, что сначала вершины с числами \pseudocode{10} и \pseudocode{2} были поменяны местами):\\
    \begin{tikzpicture}[
        scale=.4,
        every node/.style={scale=.8}
        ]
        \node[treenode, rectangle, minimum height=5mm, text centered] (root) at (0,0) {10};
        \node[treenode] (l) at (-4,-2) {5};
        \node[treenode] (r) at (4,-2) {4};
        \node[treenode] (ll) at (-6,-4) {8};
        \node[treenode] (lr) at (-2,-4) {6};
        \node[treenode] (rl) at (2,-4) {7};
        \node[treenode] (rr) at (6,-4) {42};
        \node[treenode] (lll) at (-7,-6) {11};
        \node[treenode] (llr) at (-5,-6) {15};
        \node[treenode] (lrl) at (-3,-6) {28};
        \node[treenode] (lrr) at (-1,-6) {9};
        \node[treenode] (rll) at (1,-6) {13};
        
        \draw[->, red] (root) -- (l);
        \draw[->, red] (root) -- (r);
        \draw[->] (l) -- (ll);
        \draw[->] (l) -- (lr);
        \draw[->] (r) -- (rl);
        \draw[->] (r) -- (rr);
        \draw[->] (ll) -- (lll);
        \draw[->] (ll) -- (llr);
        \draw[->] (lr) -- (lrl);
        \draw[->] (lr) -- (lrr);
        \draw[->] (rl) -- (rll);
    \end{tikzpicture}
    \begin{tikzpicture}[
        scale=.4,
        every node/.style={scale=.8}
        ]
        \node[treenode] (root) at (0,0) {4};
        \node[treenode] (l) at (-4,-2) {5};
        \node[treenode, rectangle, minimum height=5mm, text centered] (r) at (4,-2) {10};
        \node[treenode] (ll) at (-6,-4) {8};
        \node[treenode] (lr) at (-2,-4) {6};
        \node[treenode] (rl) at (2,-4) {7};
        \node[treenode] (rr) at (6,-4) {42};
        \node[treenode] (lll) at (-7,-6) {11};
        \node[treenode] (llr) at (-5,-6) {15};
        \node[treenode] (lrl) at (-3,-6) {28};
        \node[treenode] (lrr) at (-1,-6) {9};
        \node[treenode] (rll) at (1,-6) {13};
        
        \draw[->] (root) -- (l);
        \draw[->] (root) -- (r);
        \draw[->] (l) -- (ll);
        \draw[->] (l) -- (lr);
        \draw[->, red] (r) -- (rl);
        \draw[->, darkgreen] (r) -- (rr);
        \draw[->] (ll) -- (lll);
        \draw[->] (ll) -- (llr);
        \draw[->] (lr) -- (lrl);
        \draw[->] (lr) -- (lrr);
        \draw[->] (rl) -- (rll);
    \end{tikzpicture}
    \begin{tikzpicture}[
        scale=.4,
        every node/.style={scale=.8}
        ]
        \node[treenode] (root) at (0,0) {4};
        \node[treenode] (l) at (-4,-2) {5};
        \node[treenode] (r) at (4,-2) {7};
        \node[treenode] (ll) at (-6,-4) {8};
        \node[treenode] (lr) at (-2,-4) {6};
        \node[treenode, rectangle, minimum height=5mm, text centered] (rl) at (2,-4) {10};
        \node[treenode] (rr) at (6,-4) {42};
        \node[treenode] (lll) at (-7,-6) {11};
        \node[treenode] (llr) at (-5,-6) {15};
        \node[treenode] (lrl) at (-3,-6) {28};
        \node[treenode] (lrr) at (-1,-6) {9};
        \node[treenode] (rll) at (1,-6) {13};
        
        \draw[->] (root) -- (l);
        \draw[->] (root) -- (r);
        \draw[->] (l) -- (ll);
        \draw[->] (l) -- (lr);
        \draw[->] (r) -- (rl);
        \draw[->] (r) -- (rr);
        \draw[->] (ll) -- (lll);
        \draw[->] (ll) -- (llr);
        \draw[->] (lr) -- (lrl);
        \draw[->] (lr) -- (lrr);
        \draw[->, darkgreen] (rl) -- (rll);
    \end{tikzpicture}\\
    Нам обещали алгоритмы с помощью структур данных, вот, этой структурой данных можно сортировать. Как сортировать кучей? Очень просто. Берёшь массив, делаешь из него кучу ($n$ раз вызывая \pseudocode{insert}). Потом вытаскиваешь элементы из кучи (вызывая $n$ раз \pseudocode{get_min}+\pseudocode{remove_min}). Сложность --- $O(n\log n)$. Есть проблема. Нужен ещё один массив. Его заводить впадлу. Как без помощи второго массива превратить первый в кучу? Да на изи на самом деле. В каждый момент времени у нас первые \pseudocode{i} элементов --- это куча, а оставшиеся --- кусок нужного нам массива. То есть мы просто просеиваем вверх каждый элемент по порядку, а потом меняем местами первый элемент с последним, делаем вид, что последний элемент не при делах, корень просеиваем вниз. Массив, к сожалению, по убыванию отсортируется, но мы можем просто изменить строение кучи с $\leqslant$ на $\geqslant$.\\
    А давайте теперь покажем, что наша программа --- это и $\varOmega(n\log n)$ (в худшем случае). Давайте докажем, что \textit{кучификация} (то есть первый этап сортировки) $\varOmega(n\log n)$. По-хорошему это $\log_21+\log_22+\ldots+\log_2n$. На первую половину этого мы забьём, получаем $\log_2\frac n2+\ldots+\log_2n$. Это не меньше, чем $\frac n2\log_2\frac n2=\varOmega(n\log n)$.\\
    Оптимизация. Давайте посмотрим на следующее. Когда мы построили кучу из половины массива, оставшаяся половина будет строиться долго, потому что куча уже будет иметь высоту $\log_2n-1$, а нам в неё надо впихнуть (просеиванием вверх) $\frac n2$. Это будет долго. Вместо этого можно сделать наоборот. Можно строить кучу с конца. Тогда размер кучи достигнет $\log_2n-1$ только тогда, когда мы дойдём до последнего элемента. Точнее, нет, технически там будет два поддерева кучи, каждое $\log_2n-1$, первое --- с корнем в \pseudocode{1}, второе --- в \pseudocode{2}. То есть у нас будет всего один элемент, который мы будем просеивать на куче высотой $\log_2n$ вместо целых $\frac n2$ таковых.
    % Возможно, картинка с просеиванием вниз.
    Полученная штука работает за $O(n)$. Почему? Заметим следующее. С первой половиной элементов мы не делаем ничего. В следующей четверти находятся вершины, детьми которых являются вершины первой половины, а значит, нужно просеять эту четверть вниз, по дереву высоты 2, совершив одно действие. Следующая восьмая часть просеивается по дереву высоты 3, то есть с двумя действиями. Следующая шестнадцатая --- три действия и так далее. Чтобы доказать, что эта сумма --- $O(n)$, есть два пути. Математический выглядит так: $\sum\limits_{i=0}^{\log_2n}i\frac n{2^{i+1}}=\frac n4\sum\limits_{i=0}^{\log_2n}i\frac1{2^{i-1}}$. Обозначим $\frac12$ за $x$. Получим $\frac n4\sum\limits_{i=0}^{\log_2n}ix^{i-1}=\frac n4\sum\limits_{i=0}^{\log_2n}\frac{\mathrm d}{\mathrm dx}x^i=\frac n4\frac{\mathrm d}{\mathrm dx}\sum\limits_{i=0}^{\log_2n}x^i=\frac n4\frac{\mathrm d}{\mathrm dx}\frac{x^{\log_2n+1}-1}{x-1}=\frac n4\frac{\log_2nx^{\log_2n+1}-(\log_2n+1)x^{\log_2n}+1}{(x-1)^2}$. Вспоминаем, что $x=\frac12$. $\frac n4\frac{\log_2n\frac12\frac1n-(\log_2n+1)\frac1n+1}{(\frac12-1)^2}=n(\log_2n\frac12\frac1n-(\log_2n+1)\frac1n+1)=n-1-\frac12\log_2n=O(n)$. Но это сложно и как-то нелегально выглядит. Поэтому, вот способ проще. Сначала мы сложим $\frac n4$, $\frac n8$, $\frac n{16}$ и так далее. Получим $\frac n2$. Тут у нас закончатся все слагаемые вида $\frac n4$, а всех остальных останется на 1 меньше, чем было. Потом сложим $\frac n8$, $\frac n{16}$, $\frac n{32}$ и так далее. Получим $\frac n4$. Тут у нас закончатся слагаемые вида $\frac n8$, а всех остальных станет уже на 2 меньше. Делаем так дальше, получаем сумму $\frac n2+\frac n4+\frac n8+\cdots$. Эта сумма равна $n$.\\
    Возникает вопрос, почему получились разные ответы? А на самом деле потому, что во втором случае есть погрешность. Например, при $n=16$ мы нигде и ни с чем не будем складывать $\frac n{32}$. А в указанном способе мы складываем, считая что указанные ряды бесконечны. Тем не менее, мы получили функцию, б\'{о}льшую, чем нужно, что докажет нам $O(n)$, но не докажет $\varOmega(n)$, в отличие от первого способа.
    \paragraph{\undercolorblack{orange}{Рандомизованные алгоритмы}.}
    Как обычно работает алгоритм? Ему дают ввод, он выводит вывод. У рандомизованных есть ещё один вход --- случайные числа. На первый взгляд странно, зачем это. Но в некоторых алгоритмах бывают полезны хорошие генераторы случайных чисел. И сегодня мы напишем сортировку. Она будет за $O(n\log n)$, то есть никаких откровений, но иногда действительно бывает, когда рандомизованные алгоритмы работают асимптотически быстрее.\\
    Как оценивается время работы случайного алгоритма? У него же зависит время от случайных чисел, а значит само является случайной величиной. И вместо того, чтобы оценивать время $T(n)$, мы будем оценивать его математическое ожидание $T^*(n)=\mathbb E(T(n))$. Заметим, что, несмотря на то, что мы хотим худшее время, мы не хотим худший рандом, а хотим худшее мат. ожидание. Проблема. Мы не знаем, что такое мат. ожидание, у нас не было теории вероятности. Но и пофиг, функции у нас будут тупыми, справимся.
    \subparagraph{\undercolorblack{orange}{Быстрая сортировка (a.k.a. сортировка Хоара)}.}
    Есть массив. Очень сильно думаем, что все элементы разные. Выберем случайный элемент массива (т.е. возьмём случайное число \pseudocode{i} от \pseudocode{0} до \pseudocode{n-1} и выберем элемент \pseudocode{a[i]} с номером \pseudocode{i}). Дальше мы возьмём, поместим все элементы, меньшие \pseudocode{a[i]} левее него, а б\'{о}льшие либо равные --- правее. Рекурсивно сортируем обе части.
    \begin{minted}{\pseudo}
        sort(a):
            x = a[rand(0..n - 1)]
            a1, a2 = split(a, x)
            sort(a1)
            sort(a2)
            a = a1 + a2
    \end{minted}
    Непонятно только, что делать с функцией \pseudocode{split}. За какое время мы можем поделить массив на часть, б\'{о}льшую данной, и часть, меньшую данной? За $O(n)$. О'кей, а что с рекурсией. Пусть \pseudocode{len(a1)} --- это \pseudocode{n1}, a \pseudocode{len(a2)} --- это \pseudocode{n2}. То есть $T(n)=n+T(n1)+T(n2)$. Почему это $O(n\log n)$? Ну, пока не понятно. Что такое мат. ожидание? Это сумма значений, умноженных на их вероятности. $\mathbb E(T)=\sum_tt\mathbb P(T=t)$. Будем считать, что, если \pseudocode{x} из средней трети массива, то всё хорошо, а иначе --- не хорошо. Распишем мат. ожидание так. $T^*(n)\leqslant n+\frac13\left(T^*(\frac n3)+T^*(\frac{2n}3)\right)+\frac23T^*(n)$. Почему так? Ну, со второй частью (которая при $\frac23$) понятно, мы не можем поделить массив хуже (то есть для большего времени работы), чем поделить его по минимальному элементу. Ну, и с первой частью то же самое, самый плохой случай --- когда мы на следующем шаге делим нашу центральную треть по её минимуму.\\
    Итак, $T^*(n)\leqslant n+\frac13\left(T^*(\frac n3)+T^*(\frac{2n}3)\right)+\frac23T^*(n)$. Нам хочется доказать, что $T^*\leqslant cn\log_3n$. Докажем это по индукции. Остаётся $T^*(n)\leqslant n+\frac13\left(c\frac n3\log_3\frac n3+c\frac{2n}3\log_3\frac{2n}3\right)+\frac23nT^*(n)$. Вычтем последнее слагаемое из обеих частей, после чего всё умножим всё на 3. Получим $T^*(n)\leqslant3n+c\frac n3\log_3\frac n3+c\frac{2n}3\log_3\frac{2n}3=3n+c\frac n3(\log_3n-1)+c\frac{2n}3(\log_3n-1+\log_32)=cn\log_3n+3n-c\left(\frac n3+\frac{2n}3-\frac{2n}3\log_32\right)=cn\log_3n+n\left(3-c\left(1-\frac23\log_32\right)\right)$. Заметим, что $c>\frac3{1-\frac23\log_32}>0$ подходит.\\
    Давайте напишем это, но умнее. В \pseudocode{sort} мы будем передавать кусок массива. А функция \pseudocode{split} будет запускаться от куска массива и \pseudocode{x}, внутри куска строя нужную нам структуру. А ещё возвращать то, где должен лежать \pseudocode{x}.
    \begin{minted}{\pseudo}
        sort(l, r):
            if r - l <= 1
                return # Already sorted.
            x = a[rand(l, r - 1)] # `rand` returns value in closed interval.
            m = split(l, r, x)
            sort(l, m)
            sort(m, r)

        split(l, r, x):
            m = l
            for i in [l..r]
                if a[i] < x
                    swap(a[i], a[m++])
            return m
    \end{minted}
    Иллюстрация для \pseudocode{split}. Если элемент с номером \pseudocode{i} больше либо равен \pseudocode{x} нам абсолютно незачем его трогать. Если же меньше, то запихнём его в конец первого блока и увеличим индекс \pseudocode{m}:\\
    \begin{tikzpicture}[scale=.45]
        \draw[step=1cm,black,very thin,pattern=north west lines,pattern color=darkgreen] (0,0) rectangle (5,1);
        \draw[step=1cm,black,very thin,pattern=north east lines,pattern color=darkgreen] (5,0) rectangle (9,1);
        \draw[step=1cm,black,very thin] (11.5,0) -- (9,0) -- (9,1) -- (11.5,1);

        \filldraw[color=black,fill=white,thin] (5.5,.5) circle (.25);
        \filldraw[color=black,fill=white,thin] (9.5,.5) circle (.25);

        \node at (5,1.25) {$m$};
        \node at (9.5,1.25) {$i$};
        \node[scale=.75] at (3,-.25) {$<x$};
        \node[scale=.75] at (7.5,-.25) {$\geqslant x$};

        \draw[<->, thin, dashed] (5.5,-.125) arc (240:300:4);
    \end{tikzpicture}\\
    Жрёт ли это память? Только для стека рекурсии. То есть нам нужно хранить, какими были \pseudocode{l}, \pseudocode{r}, \pseudocode{x} и \pseudocode{m} на каждом шаге, чтобы после выхода из рекурсивной функции восстановить эти значения. И мат. ожидание количества этой памяти, понятно, $O(\log n)$. Про стеки разговор пойдёт боле детально позже.
    \subparagraph{\undercolorblack{orange}{$k$-порядковая статистика}.}
    Есть массив. Хочется понять, какой элемент будет на $k$-том месте, если массив отсортировать. Именно этот элемент и называется $k$-порядковой статистикой. За $O(n\log n)$ понятно, как решать. Но можно решить конкретно эту задачу быстрее.
    \subparagraph{\undercolorblack{orange}{Алгоритм Хуара}.}
    Возьмём массив. Возьмём его случайных элемент. Поделим всё на два куска, такие же, как и в сортировке. $k$-тый элемент лежит либо в левой части, либо в правой. Если $k$-тый элемент лежит в левой части, правую нам можно вообще игнорировать. И наоборот. Дальше снова делаем рекурсию:
    \begin{minted}{\pseudo}
        find(l, r, k): # `k` is GLOBAL index in array.
            if r - l <= 1
                return l
            x = a[rand(l, r - 1)]
            m = split(l, r, x) # `m` is also global, as we can see.
            if k < m
                return find(l, m)
            else
                return find(m, r)
    \end{minted}
    Тут с оценкой асимптотики всё точно также. $T^*(n)=n+\frac13T^*\left(\frac{2n}3\right)+\frac23T^*(n)$, а значит $T^*(n)\leqslant3n+c\frac{2n}3$. Мы хотим, чтобы выполнялось следующее: $3n+c\frac{2n}3<cn$. Нам хватит $c\geqslant9$.\\
    А что с детерминизмом? Обычно, когда находят рандомизированный алгоритм, пытаются найти детерминированный. Потому что, во-первых, может быть что-то критически важное, зависящее от скорости алгоритма, даже если ваш ядерный реактор будет взрываться раз в 20 лет, это всё равно не хорошо. А во-вторых, рандом непонятно, откуда брать. Надо как-то гарантировать, что это именно чистый, хороший рандом, а как он в компьютере реализован, вообще хз.
    \subparagraph{\undercolorblack{orange}{Алгоритм Блюма-Флойда-Пратта-Ривеста-Тарьяна}.}
    Идея состоит в том, чтобы делать всё то же самое, но \pseudocode{x} искать не случайно. И у нас на это линейное время. Поделим наш массив на блоки по 5 элементов. Каждый блок размера \pseudocode{5} сортируется за $O(1)$. В каждом блоке возьмём средний элемент. Среди них найдём медиану. Медиана --- это что? Медиана --- это $\frac n2$-порядковая статистика. Рекурсивно найдём её. Утверждается, что эта самая медиана --- это хороший \pseudocode{x} для разделения и рекурсивного запуска.\\
    Асимптотика. Посмотрим на медиану медиан (\pseudocode{x}). У неё половина элементов меньше неё. А это значит, что в тех пятёрках, медиана которых меньше \pseudocode{x} есть по 3 элемента, меньших \pseudocode{x}. То есть всего у нас есть не более $\frac{3n}{10}$ элементов, от которых мы потом запустим рекурсию. То есть $T(n)=n+T\left(\frac n5\right)+T\left(\frac{7n}{10}\right)$. $T(n)\leqslant n+c\frac{9n}{10}$. Тут $n\geqslant10$.
    \paragraph{\undercolorblack{orange}{Оптимальная асимптотика сортировки}.}
    Доказать, что нельзя что-то делать быстрее, что за какое-то время, обычно очень сложно. Это делается как-нибудь так: Докажем, что нельзя найти минимум меньше, чем за линию. Пусть можно. Тогда мы не посмотрели на хотя бы один элемент. А он мог оказаться максимумом. Проблема.\\
    Сортировка, основанная на сравнениях (то есть такая, в которой мы смотрим на массив только для сравнений), не может работать быстрее, чем за $O(n\log n)$. Докажем это. Построим следующее дерево. У нас на каждом шаге алгоритм что-то сравнивает и в зависимости от результата алгоритм делает либо что-то одно, либо что-то другое. В конце алгоритм даёт ответ. Сколько у нас есть возможных ответов? Ну, не менее, чем $n!$, потому что у массива любая перестановка может быть его отсортированной версией. То есть листов у этого дерева сравнений хотя бы $n!$. Какова высота этого дерева? Ну, $\log_2(n!)=\sum\limits_{i=0}^n\log_2i\geqslant\frac n2\log_2\frac n2=\varOmega(n\log n)$. Это доказательство полностью ломается, если мы вдруг умеем делать что-то ещё, ведь мы при помощи этого сможем отсортировать быстрее.
    \paragraph{\undercolorblack{orange}{Сортировки, быстрее $n\log n$}.}
    Пусть в массиве \pseudocode{a[0..n-1]} лежат маленькие (меньшие \pseudocode{m}) числа. Чтобы таковой отсортировать можно посчитать количество каждого элемента. Это сортировка подсчётом.
    \begin{minted}{\pseudo}
        # `cnt` is m-length array filled with zeros.
        for i in [0:n-1]
            cnt[a[i]]++
    \end{minted}
    Но у нас есть мем. Это не то, чтобы сортировка. Если мы сортируем, например, пары по ключу, то ключи-то у нас сохранятся, а значения --- нет, ведь мы по факту забываем массив, а потом создаём новый. Ну, например, пусть у нас массив \pseudocode{a} выглядит как \pseudocode{[<2, A>, <1, B>, <1, C>, <0, D>, <2, E>, <2, F>, <0, G>]}. Тогда модифицируем алгоритм так. \pseudocode{cnt} --- это не массив чиселок, а массив списков. И мы кладём во второй список пару \pseudocode{<2, A>}, в первый --- \pseudocode{<1, B>} и т.д. Если у нас правильно реализованы списки, то сортировка работает за линейное время. А без списков как сделать? Можно хранить обычные количества элементов и сказать, что (в нашем случае) первые 2 места зарезервированы под ключи \pseudocode{0}, вторые два --- под \pseudocode{1}, последние три --- под \pseudocode{2}. Эти зарезервированные куски массива назовём корзинками. Как быстро добавлять элемент в нужное место корзинки? Просто хранить позиции. Назовём новый массив \pseudocode{p}. Изначально он является префиксной суммой \pseudocode{cnt}, а дальше делаем так:
    \begin{minted}{\pseudo}
        for i in [0..n-1]
            x = a[i].key
            b[p[x]++] = a[i]
    \end{minted}
    Время у нас получается $O(n+m)$, памяти --- тоже.\\
    Усложним задачу. Пусть у нас количество все числа не от \pseudocode0 до \pseudocode{m-1}, а от \pseudocode0 до \pseudocode{m**2-1}. Каждое число \pseudocode{x} мы разобьём на две компоненты --- \mintinline{\pseudo}{x \% m} и \pseudocode{x / m}.
    Теперь мы просто сортируем числа во возрастанию второго, а если вторые равны --- то первого. Как это сделать? Да очень просто. Наша сортировка стабильна, поэтому мы можем отсортировать сначала по остатку, а потом по целой части. Да, именно в таком порядке. После первой операции у нас будет все элементы с меньшим остатком раньше элементов больше, а после второй числа с равными первыми цифрами останутся отсортированы по второй.
    \subparagraph{\undercolorblack{orange}{Цифровая сортировка}.}
    Давайте просто обобщим эту штуку на произвольную степень \pseudocode{m}. Можно взять, например, \pseudocode{m} равное \pseudocode{10}, а потом сортировкой подсчётом сортировать сначала по самой младшей цифре, потом по следующей и так далее до самой старшей. Итого у нас получится $O((m+n)\cdot k)$, где $k$ --- максимальное количество цифр. То есть оптимально брать \pseudocode{m} примерно равное \pseudocode{n}.
    \paragraph{\undercolorblack{orange}{Сортировочная сеть}.}
    Пусть у нас есть только следующая операция:
    \begin{minted}{\pseudo}
        f(i, j):
            if a[i] > a[j]
                swap(a[i], a[j])
    \end{minted}
    Попробуем ей сортировать. Например, пузырёк работает так:
    \begin{minted}{\pseudo}
        for i in [n-1..0]
            for j in [0..i-1]
                f(j, j+1)
    \end{minted}
    Это можно изобразить так:\\
    \begin{tikzpicture}
        \draw[thin] (-4,-2) -- (3,-2);
        \draw[thin] (-4,-1) -- (3,-1);
        \draw[thin] (-4,0) -- (3,0);
        \draw[thin] (-4,1) -- (3,1);
        \draw[thin] (-4,2) -- (3,2);
        \draw[->,thick] (-3.5,2) -- (-3.5,1);
        \draw[->,thick] (-3,1) -- (-3,0);
        \draw[->,thick] (-2.5,0) -- (-2.5,-1);
        \draw[->,thick] (-2,-1) -- (-2,-2);

        \draw[->,thick] (-1,2) -- (-1,1);
        \draw[->,thick] (-.5,1) -- (-.5,0);
        \draw[->,thick] (0,0) -- (0,-1);

        \draw[->,thick] (1,2) -- (1,1);
        \draw[->,thick] (1.5,1) -- (1.5,0);

        \draw[->,thick] (2.5,2) -- (2.5,1);
    \end{tikzpicture}
    Это и называется сортировочной сетью, и у неё есть 2 параметра. Первый --- это количество таких сравнений. Но заметим, что у нас можно одновременно выполнять некоторые операции, и второй --- это и есть количество операций, при условии, что мы выполняем одновременно всё, что можем. У нас, например, $O(n^2)$ операций всего и $O(n)$, если выполнять параллельно.
    \subparagraph{\undercolorblack{orange}{Битонная сортировка}.}
    Рассмотрим последовательность, которая сначала возрастает, а потом убывает. Любой её циклический сдвиг называется битонной последовательностью. Битонная сортировка сортирует такие последовательности. При наличии мощной параллельности.\\
    Теорема. Сортирующая сеть сортирует массив тогда и только тогда, когда она сортирует любой массив из нулей и единиц. Следствие вправо очевидно. Следствие влево. Рассмотрим наименьший элемент. Вместо того массива, который мы хотим отсортировать, впихнём вместо наименьшего элемента ноль, а вместо всех остальных --- единицу. А значит наш наименьший элемент будет идти по тем же путям, что и этот самый ноль. А значит после сортировки наименьший элемент будет на правильном месте. Окей, как доказать, что второй элемент будет на своём месте? Для этого рассмотрим другой массив из нулей и единиц, вместо двух наименьших элементов поставим 0, а для остальных --- 1. И так далее. Доказано.\\
    Как выглядит битонная последовательность из нулей и единиц. Ну, плохо. У неё есть нули в начале и конце и единицы в середине, или наоборот. Как это сортировать?
    Пусть количество элементов --- степень двойки. Тогда сделаем так:\\
    \begin{tikzpicture}
        \foreach\x in {0,...,7} {
            \draw[thin] (0,\x) -- (4,\x);
        }
        \draw[thick,->] (.5,7) -- (.5,3);
        \draw[thick,->] (.75,6) -- (.75,2);
        \draw[thick,->] (1,5) -- (1,1);
        \draw[thick,->] (1.25,4) -- (1.25,0);
    \end{tikzpicture}\\
    Утверждение. Это превратит нашу битонную последовательность в две других. Притом, всё слева меньше всего справа. Докажем это перебором случаев.\\
    \begin{enumerate}[1.]
        \item Если все единички только слева.
        \begin{tikzpicture}[scale=.25, every node/.style={scale=0.55}]
            \draw[step=1cm,black,very thin] (0,0) grid (16,1);
            \node at (.5,.5) {0};
            \node at (1.5,.5) {1};
            \node at (2.5,.5) {1};
            \node at (3.5,.5) {1};
            \node at (4.5,.5) {0};
            \node at (5.5,.5) {0};
            \node at (6.5,.5) {0};
            \node at (7.5,.5) {0};
            \node at (8.5,.5) {0};
            \node at (9.5,.5) {0};
            \node at (10.5,.5) {0};
            \node at (11.5,.5) {0};
            \node at (12.5,.5) {0};
            \node at (13.5,.5) {0};
            \node at (14.5,.5) {0};
            \node at (15.5,.5) {0};
            
            \draw[dashed,thick,red] (8,-1) -- (8,2);
        \end{tikzpicture}
        превращается в
        \begin{tikzpicture}[scale=.25, every node/.style={scale=0.55}]
            \draw[step=1cm,black,very thin] (0,0) grid (16,1);
            \node at (.5,.5) {0};
            \node at (1.5,.5) {0};
            \node at (2.5,.5) {0};
            \node at (3.5,.5) {0};
            \node at (4.5,.5) {0};
            \node at (5.5,.5) {0};
            \node at (6.5,.5) {0};
            \node at (7.5,.5) {0};
            \node at (8.5,.5) {0};
            \node at (9.5,.5) {1};
            \node at (10.5,.5) {1};
            \node at (11.5,.5) {1};
            \node at (12.5,.5) {0};
            \node at (13.5,.5) {0};
            \node at (14.5,.5) {0};
            \node at (15.5,.5) {0};
            
            \draw[dashed,thick,red] (8,-1) -- (8,2);
        \end{tikzpicture}.
        \item Если все единички справа. \begin{tikzpicture}[scale=.25, every node/.style={scale=0.55}]
            \draw[step=1cm,black,very thin] (0,0) grid (16,1);
            \node at (.5,.5) {0};
            \node at (1.5,.5) {0};
            \node at (2.5,.5) {0};
            \node at (3.5,.5) {0};
            \node at (4.5,.5) {0};
            \node at (5.5,.5) {0};
            \node at (6.5,.5) {0};
            \node at (7.5,.5) {0};
            \node at (8.5,.5) {0};
            \node at (9.5,.5) {1};
            \node at (10.5,.5) {1};
            \node at (11.5,.5) {1};
            \node at (12.5,.5) {0};
            \node at (13.5,.5) {0};
            \node at (14.5,.5) {0};
            \node at (15.5,.5) {0};
            
            \draw[dashed,thick,red] (8,-1) -- (8,2);
        \end{tikzpicture}
        остаётся как есть.
        \item Если единички пересекают линию разделения, и их мало. \begin{tikzpicture}[scale=.25, every node/.style={scale=0.55}]
            \draw[step=1cm,black,very thin] (0,0) grid (16,1);
            \node at (.5,.5) {0};
            \node at (1.5,.5) {0};
            \node at (2.5,.5) {0};
            \node at (3.5,.5) {0};
            \node at (4.5,.5) {0};
            \node at (5.5,.5) {0};
            \node at (6.5,.5) {0};
            \node at (7.5,.5) {1};
            \node at (8.5,.5) {1};
            \node at (9.5,.5) {1};
            \node at (10.5,.5) {0};
            \node at (11.5,.5) {0};
            \node at (12.5,.5) {0};
            \node at (13.5,.5) {0};
            \node at (14.5,.5) {0};
            \node at (15.5,.5) {0};
            
            \draw[dashed,thick,red] (8,-1) -- (8,2);
        \end{tikzpicture}
        превратится в
        \begin{tikzpicture}[scale=.25, every node/.style={scale=0.55}]
            \draw[step=1cm,black,very thin] (0,0) grid (16,1);
            \node at (.5,.5) {0};
            \node at (1.5,.5) {0};
            \node at (2.5,.5) {0};
            \node at (3.5,.5) {0};
            \node at (4.5,.5) {0};
            \node at (5.5,.5) {0};
            \node at (6.5,.5) {0};
            \node at (7.5,.5) {0};
            \node at (8.5,.5) {0};
            \node at (9.5,.5) {1};
            \node at (10.5,.5) {0};
            \node at (11.5,.5) {0};
            \node at (12.5,.5) {0};
            \node at (13.5,.5) {0};
            \node at (14.5,.5) {1};
            \node at (15.5,.5) {1};
            
            \draw[dashed,thick,red] (8,-1) -- (8,2);
        \end{tikzpicture}
        \item Если единички пересекают линию разделения, и их много (т.е. существует место, где мы сравниваем две единички).
        \begin{tikzpicture}[scale=.25, every node/.style={scale=0.55}]
            \draw[step=1cm,black,very thin] (0,0) grid (16,1);
            \node at (.5,.5) {0};
            \node at (1.5,.5) {0};
            \node at (2.5,.5) {0};
            \node at (3.5,.5) {0};
            \node at (4.5,.5) {1};
            \node at (5.5,.5) {1};
            \node at (6.5,.5) {1};
            \node at (7.5,.5) {1};
            \node at (8.5,.5) {1};
            \node at (9.5,.5) {1};
            \node at (10.5,.5) {1};
            \node at (11.5,.5) {1};
            \node at (12.5,.5) {1};
            \node at (13.5,.5) {1};
            \node at (14.5,.5) {0};
            \node at (15.5,.5) {0};
            
            \draw[dashed,thick,red] (8,-1) -- (8,2);
        \end{tikzpicture}
        превратится в
        \begin{tikzpicture}[scale=.25, every node/.style={scale=0.55}]
            \draw[step=1cm,black,very thin] (0,0) grid (16,1);
            \node at (.5,.5) {0};
            \node at (1.5,.5) {0};
            \node at (2.5,.5) {0};
            \node at (3.5,.5) {0};
            \node at (4.5,.5) {1};
            \node at (5.5,.5) {1};
            \node at (6.5,.5) {0};
            \node at (7.5,.5) {0};
            \node at (8.5,.5) {1};
            \node at (9.5,.5) {1};
            \node at (10.5,.5) {1};
            \node at (11.5,.5) {1};
            \node at (12.5,.5) {1};
            \node at (13.5,.5) {1};
            \node at (14.5,.5) {1};
            \node at (15.5,.5) {1};
            
            \draw[dashed,thick,red] (8,-1) -- (8,2);
        \end{tikzpicture}
    \end{enumerate}
    Если же у нас в центре нолики, а вокруг единички, а не наоборот, то происходит то же самое.\\
    А значит мы можем рекурсивно запуститься от каждой половинки, но они не зависят, поэтому это делается параллельно. Как быстро это работает? Если делать всё последовательно, получится $O(n\log n)$. А вот если параллельно, то вообще $O(\log n)$. Итого алгоритм можно проиллюстрировать так:\\
    \begin{tikzpicture}
        \foreach\x in {0,...,7} {
            \draw[thin] (0,\x) -- (4,\x);
        }
        \draw[thick,->] (.5,7) -- (.5,3);
        \draw[thick,->] (.75,6) -- (.75,2);
        \draw[thick,->] (1,5) -- (1,1);
        \draw[thick,->] (1.25,4) -- (1.25,0);
        
        \draw[thick,->] (2.25,7) -- (2.25,5);
        \draw[thick,->] (2.5,6) -- (2.5,4);
        \draw[thick,->] (2.25,3) -- (2.25,1);
        \draw[thick,->] (2.5,2) -- (2.5,0);
        
        \draw[thick,->] (3.5,7) -- (3.5,6);
        \draw[thick,->] (3.5,5) -- (3.5,4);
        \draw[thick,->] (3.5,3) -- (3.5,2);
        \draw[thick,->] (3.5,1) -- (3.5,0);
    \end{tikzpicture}\\
    Есть один нюанс. Как сделать последовательность битонной? Давайте левую часть отсортируем рекурсивно по возрастанию, а правую --- по убыванию. Получится, что у нас $O(\log n)+O(\log n-1)+O(\log n-2)+\cdots=O(\log^2n)$ параллельных операций и $O(n\log^2 n)$ последовательных. То есть итоговая сортировочная сеть для массива 8 элементов выглядит так:\\
    \begin{tikzpicture}
        \foreach\x in {0,...,7} {
            \draw[thin] (-2.75,\x) -- (4,\x);
        }
        \draw[darkgreen,thick,->] (-2.25,7) -- (-2.25,6);
        \draw[darkgreen,thick,<-] (-2.25,5) -- (-2.25,4);
        \draw[darkgreen,thick,->] (-2.25,3) -- (-2.25,2);
        \draw[darkgreen,thick,<-] (-2.25,1) -- (-2.25,0);

        \draw[black!50!red,thick,->] (-1.25,7) -- (-1.25,5);
        \draw[black!50!red,thick,->] (-1.5,6) -- (-1.5,4);
        \draw[black!50!red,thick,<-] (-1.25,3) -- (-1.25,1);
        \draw[black!50!red,thick,<-] (-1.5,2) -- (-1.5,0);

        \draw[black!50!red,thick,->] (-.5,7) -- (-.5,6);
        \draw[black!50!red,thick,->] (-.5,5) -- (-.5,4);
        \draw[black!50!red,thick,<-] (-.5,3) -- (-.5,2);
        \draw[black!50!red,thick,<-] (-.5,1) -- (-.5,0);

        \draw[blue,thick,->] (.5,7) -- (.5,3);
        \draw[blue,thick,->] (.75,6) -- (.75,2);
        \draw[blue,thick,->] (1,5) -- (1,1);
        \draw[blue,thick,->] (1.25,4) -- (1.25,0);
        
        \draw[blue,thick,->] (2.25,7) -- (2.25,5);
        \draw[blue,thick,->] (2.5,6) -- (2.5,4);
        \draw[blue,thick,->] (2.25,3) -- (2.25,1);
        \draw[blue,thick,->] (2.5,2) -- (2.5,0);
        
        \draw[blue,thick,->] (3.5,7) -- (3.5,6);
        \draw[blue,thick,->] (3.5,5) -- (3.5,4);
        \draw[blue,thick,->] (3.5,3) -- (3.5,2);
        \draw[blue,thick,->] (3.5,1) -- (3.5,0);
    \end{tikzpicture}\\
    То есть сначала у нас есть несколько последовательностей длины 2. Они битонные всегда. Потом мы сортируем их в нужном нам порядке, чтобы они образовали битонные последовательности длины 4 (это зелёные стрелки). После этого мы сортируем последовательности длины 4, чтобы получить одну битонную последовательность длины 8 (коричневые стрелки). Потом просто сортируем её (синие стрелки).
    \paragraph{\undercolorblack{orange}{Двоичный поиск}.}
    У нас есть отсортированный массив \pseudocode{a}. Например, \pseudocode{[1, 3, 4, 8, 10, 13, 17, 21]}. Мы хотим найти такое \pseudocode{i}, что \pseudocode{a[i] == x}. Сначала неверное решение. Мы ищем \pseudocode{x} в каком-то диапазоне от \pseudocode{l} до \pseudocode{r}.
    \begin{minted}{\pseudo}
        l = 0
        r = n - 1
        while r >= l
            m = (l + r) / 2
            if a[m] < x
                # All elements from `l` to `m` are less than `x` too
                l = m + 1
            else if a[m] > x
                # All elements from `m` to `r` are greater than `x`
                r = m - 1
            else
                return m
        return Ø
    \end{minted}
    Это плохой код, в нём очень много \pseudocode{if}'ов и \pseudocode{+ 1}/\pseudocode{- 1}. К тому же он очень плохо адаптируется. Обычно мы хотим искать элемент, близкий к \pseudocode{x}, а не равный \pseudocode{x}. Если мы ищем именно \pseudocode{x}, проще было бы хэш-таблицу написать. Хорошо, а что такое <<близкий к \pseudocode{x} элемент>>? Ну, например, самый левый элемент, больше либо равный \pseudocode{x}. Если таких элементов нет, то обычно возвращают \pseudocode{n} (или \mintinline{cpp}{container.end()} в C++). Хорошо, теперь правильный алгоритм. У нас будет два указателя, \pseudocode{l} и \pseudocode{r}. Мы будем поддерживать простой инвариант, \pseudocode{a[l] < x && a[r] >= x}. Всё понятно теперь. Кроме того, чему изначально равны \pseudocode{l} и \pseudocode{r}. Чтобы инвариант сразу не обламывался, \pseudocode{l = -1}, \pseudocode{r = n}. Иногда хочется добавить в массив руками в начало \mintinline{cpp}{INT_MIN}, а в конец --- \mintinline{cpp}{INT_MAX}, но совершенно никогда не нужно, у нас не будет запросов к этому элементу массива.
    \begin{minted}{\pseudo}
        l = -1
        r = n
        while r > l + 1
            m = (l + r) / 2
            if a[m] >= x
                r = m
                # Invariant is OK
            else
                l = m
                # Invariant is also OK
        return r
    \end{minted}
    А что если мы хотим искать максимальный элемент, меньший x? Тогда код изменится незначительно:
    \begin{minted}{\pseudo}
        l = -1
        r = n
        while r > l + 1
            m = (l + r) / 2
            if a[m] > x
                r = m
            else
                l = m
            return l
    \end{minted}
    \subparagraph{\undercolorblack{orange}{Баянная задачка про квадрат}.}
    У нас есть \pseudocode{n} прямоугольников одинакового размера \pseudocode{a}\texttimes\pseudocode{b}. Хочется найти такое минимальное \pseudocode{x}, что в квадрат размера \pseudocode{x}\texttimes\pseudocode{x} влезут все наши прямоугольники. Поворачивать прямоугольники нельзя. При чём тут двоичный поиск? Если для какого-то \pseudocode{x} условие выполняется, то для любого б\'{о}льшего --- тоже, а если для \pseudocode{x} условие не выполняется, то и для любого меньшего --- тоже. В какой квадрат точно нельзя впихнуть \pseudocode{n} прямоугольников? Если все числа положительные, то в пустой (\pseudocode{l = 0}) нельзя. В какой точно можно? Ну, например, в \pseudocode{r = n * max(a, b)}. Пишем код:
    \begin{minted}{\pseudo}
        l = -1
        r = n * max(a, b)
        while r > l + 1
            m = (l + r) / 2
            if good(m)
                r = m
            else
                l = m
        return l
    \end{minted}
    Осталось лишь понять, что такое \pseudocode{good(m)}. Это функция, которая определяет, можно ли впихнуть наши \pseudocode{n} прямоугольников в квадрат \pseudocode{m}\texttimes\pseudocode{m}. Это просто.
    \begin{minted}{\pseudo}
        good(x):
            return (x / a) * (x / b) >= n
    \end{minted}
    Есть лишь одна проблема. Переполнения. Пусть у нас \pseudocode{a} большое, а \pseudocode{b} не большое. Тогда \pseudocode{r} (и \pseudocode{m}) может влезть только в \mintinline{cpp}{long long}, а промежуточный результат \pseudocode{good(m)} даже туда не влезет. Как жить?
    Самый простой (и грустный) вариант --- Python. Более интересный вариант --- \mintinline{cpp}{double}. Но тут нужно очень аккуратно жить. Например, надо посчитать \pseudocode{good(m)} в \mintinline{cpp}{double}, если результат сильно больше \pseudocode{n}, то хорошо, а если не сильно, значит в \mintinline{cpp}{long long} нормально посчитается. Ещё можно вместо выражения \pseudocode{A * B >= C} писать \pseudocode{B >= (C + A - 1) / A}. Но самый правильный вариант --- искать не очень большое начальное \pseudocode{r}. Вот так вот:
    \begin{minted}{\pseudo}
        r = 1
        while !good(r)
            r *= 2
    \end{minted}
    Тогда \pseudocode{r} будет не сильно больше ответа, а значит всё должно влезть. Более того, можно один шаг поиска пропустить, присвоив в \pseudocode{l} \pseudocode{r / 2}.
    \subparagraph{\undercolorblack{orange}{Абстрактная задача на минимакс}.}
    Задача на минимакс --- это вот такая штука: найти\\
    $\min\limits_{\pseudocode{x}}\max(f_1(\pseudocode{x});f_2(\pseudocode{x});\ldots;f_n(\pseudocode{x}))=\pseudocode{A}$. Как искать \pseudocode{A}, если мы умеем проверять условие \pseudocode{A <= y}? Очень просто. $\pseudocode{A}=\min\pseudocode{y}\colon\pseudocode{A <= y}$. Пусть мы умеем отвечать на вопрос, существует ли такой \pseudocode{x}, что $\max(f_1(\pseudocode{x});f_2(\pseudocode{});\ldots;f_n(\pseudocode{x}))\leqslant\pseudocode{y}$. Тогда несложно заметить, что нам нужно лишь найти минимальный такой \pseudocode{y}, что \pseudocode{x} существует. При этом он существует начиная с какого-то момента, а значит это самое \pseudocode{y} можно найти двоичным поиском.\\
    Пример. У нас есть мужички. Они имеют координаты и скорости. Они хотят собраться в одной точке за минимальное время. Тогда время, за которое они соберутся --- это $T(x)=\max\frac{|x_i-x|}{v_i}$. Именно это нам и надо минимизировать. Тогда давайте искать \pseudocode{x} такой, что $\max\frac{|x_i-\pseudocode{x}|}{v_i}\leqslant\pseudocode{y}$. Это значит, что для любого $i$ $\frac{|x_i-\pseudocode{x}|}{v_i}\leqslant\pseudocode{y}$. У этого даже физический смысл есть, может ли $i$-тый мужик добежать за \pseudocode{y} времени до точки \pseudocode{x}. Если нам дан какой-то \pseudocode{y}, то нам нужно лишь понять, пересекаются ли $v_i$-окрестности точек $x_i$. Если пересекаются, то любую точку пересечения можно выбрать за \pseudocode{x}. Пересекать отрезки можно достаточно быстро, а значит мы можем быстро ответить на вопрос <<Могут ли мужики за время \pseudocode{y} собраться вместе?>> А значит можно провести двоичный поиск по \pseudocode{y} и найти минимальный такой \pseudocode{y}, тем самым решив задачу.
    \subparagraph{\undercolorblack{orange}{Вещественный поиск}.}
    Мы как-то знаем, что $f(l)<0$, $f(r)>0$. При этом $f$ --- неубывающая функция, а значит она имеет корень. Чтобы его найти, можно делать всё тот же двоичный поиск. Проблема только в том, какое условие ставить в \pseudocode{while}? Можно ставить \pseudocode{r - l > e}, где \pseudocode{e} --- погрешность, с которой вы хотите найти корень. У этого есть проблема. Если \pseudocode{l} и \pseudocode{r} порядка $10^9$, а погрешность --- $10^{-9}$, то точности \mintinline{cpp}{double} не хватит. Есть два способа это решить. Если мы попали в бесконечный цикл, то \pseudocode{l} и \pseudocode{r} --- два соседних \mintinline{cpp}{double}'а. А значит, присвоив в \pseudocode{m} значение \pseudocode{l + (r - l) / 2}, мы получим либо \pseudocode{l}, либо \pseudocode{r}. Можно руками это и проверить. Можно проще. Можно сказать, что мы всегда совершаем не более, например, \pseudocode{60} действий. Этого хватит, и никогда не зациклится.
    \subparagraph{\undercolorblack{orange}{Троичный поиск}.}
    Пусть у нас есть функция, которая сначала убывает, а потом возрастает. Может быть и наоборот, но мы, не умаляя общности, будем рассматривать это вариант. Мы хотим найти минимум. Самый простой способ --- двоичный поиск по производной. Если же мы не знаем производную, то уже сложнее. Можно, конечно, сделать $\frac{f(x+\varepsilon)-f(x-\varepsilon)}{2\Delta}$, но тут проблема выбора $\Delta$, надо очень много думать. Есть стабильный способ. Мы знаем, что в \pseudocode{l} функция убывает, а в \pseudocode{r} --- возрастает. Давайте воткнёмся в две промежуточные точки \pseudocode{m1} и \pseudocode{m2}. Пусть в \pseudocode{f(m1) > f(m2)}. Тогда где-то на отрезке от \pseudocode{m1} до \pseudocode{m2} она убывала. А значит можно перенести \pseudocode{l} в \pseudocode{m1}. Конец.
    \begin{minted}{\pseudo}
        l = -1
        r = n
        while r > l + 1
            m1 = (2 * l + r) / 3
            m2 = (l + 2 * r) / 3
            if f(m1) > f(m2)
                l = m1
            else
                r = m2
        return l
    \end{minted}
    Это работает за $\log_{1.5}n$. Это хуже, чем двоичный логарифм. Можно ставить \pseudocode{m1} и \pseudocode{m2} поближе друг к другу, но тогда точность поменьше получится.
    \paragraph{\undercolorblack{orange}{Стеки, очереди и их друзья}.}
    Кто такой стек? Стек --- это набор элементиков, которые можно класть сверху (push) и удалять с верха (pop). Самый простой способ реализовывать стек --- на массиве. Проблема в том, что массив должен быть бесконечным. Первые \pseudocode{n} элементов хранят элементы стека.
    \begin{minted}{\pseudo}
        push(x):
            a[n++] = x
        pop():
            return a[--n]
    \end{minted}
    Массив нужно делать такого размера, чтобы в него всегда всё влезало. Если мы, например, знаем, что в стеке не бывает больше \pseudocode{N} элементов, то именно размера \pseudocode{N} мы массив и сделаем.\\
    Что такое очередь? У нас есть два места, начало и конец. Можно запихивать элементы в конец, а удалять из начала. Чтобы реализовать очередь на массиве, придётся делать массив бесконечным в две стороны. Пусть \pseudocode{h} --- это начало, а \pseudocode{t} --- место за концом. Тогда
    \begin{minted}{\pseudo}
        add(x):
            a[t++] = x
        remove():
            return a[h++]
    \end{minted}
    Проблема, то место, где хранятся элементы, переползает вправо. Даже если очередь всегда размера не больше 1, то это не гарантирует, что всё хорошо. Это фиксится очень просто, мы оба указателя берём по модулю \pseudocode{N}, то есть когда не хватает места в конце массива, мы начинаем добавлять в начало.\\
    Как эти штуки применяются? Очередь --- обычно именно тогда, когда по жизни нужна именно она. Например, когда мы обрабатываем какие-то задачи в порядке очереди. Стек в свою очередь применяется в различных алгоритмах. Во-первых, в рекурсии. Это мы уже, на самом деле, обговаривали, когда говорили про быструю сортировку. В стек кладутся все локальные переменные, параметры функции, позиция, в которое нужно вернуться по \mintinline{cpp}{return}... Это стек вызовов. Во-вторых, в парсерах удобно использовать стеки, если не хочется использовать рекурсию. Например, пусть мы хотим взять скобочную последовательность и сказать, какая закрывающаяся скобка к какой открывающей относится. Можно класть открывающие скобки в стек, а при встрече закрывающих --- доставать открывающие из стека.\\
    Итак, давайте вернёмся к проблеме размера массива. Хочется, чтобы этот самый размер как-то подстраивался под размер стека. В C++ есть \mintinline{cpp}{std::vector}, а в Java --- \mintinline{java}{java.util.ArrayList}. Как они работают? В чём проблема, сначала надо уточнить. В том, что мы не можем попросить у операционной системы <<вон те ещё 100 байтов, которые идут после нашего массива>>, они могут быть заняты кем-то другим. Что будет, если массив увеличивать понемножку? На 1, например. В таком случае каждое добавление копирует $n$ элементов, а значит добавление $n$ элементов будет занимать $O(n^2)$ времени. А как правильно? Увеличивать массив не <<на немножко>>, а <<в немножко раз>>. В 2, например. Сколько времени будет работать добавление $n$ элементов в пустой массив? Тогда у нас некоторые добавления в массив работают за $O(1)$, а некоторые --- за $O($n$)$. Надо лишь понять, как часто происходят вторые? Они происходят в степенях двойки. То есть происходят $\log_2n$ раз. Иначе говоря, у нас сумма таких долгих добавлений --- это сумма степеней двойки до $\log_2n$-ной. А это $O(n)$. И быстрых добавлений тоже меньше, чем \pseudocode{n} штук. Мы только что поняли, что \textit{в среднем} добавление одного элемента работает за $O(1)$. Что значит это \textit{в среднем}? А то, что амортизационное время работы равно $O(1)$. У этого есть формальное определение. Есть структура данных. Мы совершаем с ней какое-то количество действий. Тогда амортизированная оценка считается корректной, если сумма амортизированных времён не меньше, чем сумма настоящих времён. Давайте теперь чётко и формально докажем по определению амортизированное время работы. Суммарное время работы наших $n$ добавлений явно не больше, чем $5n$. А значит можно взять амортизированное время, равное 5, что даст нам доказательство.\\
    Зачем мы придумываем амортизированные структуры данных? Первый ответ --- потому что мы слабы, и знаем амортизированные СД, а превращать их в неамортизированные не умеем. Как с рандомизированными алгоритмами. Второй ответ --- структура данных внутри какого-то алгоритма имеет какой-то жизненный цикл, а значит нам интересно то и только то, сколько операций \textbf{суммарно} мы потратим, тогда амортизированные структуры --- ок. Если у нас есть атомная станция, который плохо от того, что один раз в очень редко датчики будут отвечать за 5 секунд, а не 2 миллисекунды, то это будет очень плохо, а в алгоритмах --- не очень.\\
    Как доказывать амортизированное время? Так, как сделали мы, мы запаримся для более сложных структур. Поэтому есть метод потенциалов (потанцевалов). Мы введём некоторую величину (потенциал), которая как-то отражает состояние структуры данных. Он не вылезает из СД сам, его надо взять из астрала самому. Он удовлетворяет некоторым свойствам.
    \begin{itemize}
        \item Начальный потенциал не больше константы (обычно вообще равен 0, у нас будет именно так).
        \item Потенциал всегда неотрицательный (на самом деле можно разрешать ему быть немножко отрицательным, но это нужно аккуратно контролировать, поэтому мы не будем).
    \end{itemize}
    Тогда амортизированное время будет отличаться от настоящего на разность потенциалов. Докажем же то, что это поможет. Если мы сложим все наши $n$ операций, то получим, что разность между суммой амортизированных времён и суммой истинных времён будет отличаться в разность начального потенциала и конечного потенциала. А начальный --- 0. Итак, пример. Когда мы делаем долгую операцию? Когда заканчивается место в массиве. В этом месте истинное время --- это $n$, амортизированное --- константа, а значит разность потенциалов должна быть где-то $-n$. То есть у нас потенциал заполненного массива должен быть большим, а полупустого --- маленьким. Тогда за потенциал можно брать разность между количеством заполненных ячеек и количеством свободных. Тогда потенциал и есть $-n$. Если же мы точно считаем все времена, и получаем, что истинные времена --- это, скажем, $5n$, то мы просто увеличиваем в 5 раз потенциал и чиллим. Если ввести такое амортизационное время, то мы как раз и обнаружим, что амортизационное время --- константа.\\
    Второй метод --- метод бухгалтерского учёта. Делает то же самое, но является более визуальным и простым для понимания. У нас есть две операции. Первая будет тратить $t$ времени и запасать $t$ времени на будущее, а вторая --- доставать время из запаса. То есть они обе по жизни тратят 0 времени, но первая имеет амортизационное время 1, а вторая --- -1. Если у нас первых операций в каждый момент времени больше либо равно вторых, то всё валидно. Пример для увеличивающегося массива. Давайте при добавлении элемента класть по две монетки в банк времени. Когда нам нужно расширить массив, мы возьмём наши запасённые монетки и тратим их на копирование. При этом заметим, что после копирования на старых элементов монеток лежать не будет. Потом мы добавим в массив ещё столько же элементов, сколько там было, положив вдвое больше монеток (то есть количество, равное количеству элементов), а значит нам как раз хватит их на копирование.\\
    Пусть вам по наследству от дедушки достались два стека. А вы очень хотите очередь. Тогда всё, что вам нужно, --- это склейка двух стеков низами. Добавление производится в первый, а удаление --- из второго. Но второй когда-нибудь закончится. Не беда! Когда закончится, переложим все элементы второго стека в первый.
    \begin{minted}{\pseudo}
        add(x):
            s1.push(x)
        remove():
            if empty(s2)
                while !empty(s1)
                    s2.push(s1.pop())
            return s2.pop()
    \end{minted}
    Докажем, что амортизированное время работы этого ужаса --- $O(1)$. Посмотрим на элемент. Его сначала положили в первый стек, потом переложили во второй, потом достали. То есть для каждого элемента мы в любом случае совершаем ровно 4 действия. А элементов у нас $n$. $4n=O(n)$.
    Давайте то же самое, но с потенциалом. Какое действие выполняется долго? Перекладывание. Какой нужно сделать потенциал, чтобы он уменьшался на $n$ при перекладывании? Например, количество элементов в \pseudocode{s1}. И с монетками тоже всё понятно, мы каждому элементу при добавлении даём монетку <<на проезд>>. И тратим эти монетки при перекладывании.\\
    Мы уже упоминали, что некоторые амортизированные структуры данных можно сделать без амортизации, по-честному. Например, сделаем так с расширяющимся массивом. Пусть у нас есть полузаполненный массив. Давайте заранее заведём массив вдвое больше и при каждой вставке нового элемента будем класть два старых в заведённый массив. Тогда, когда наш маленький массив заполнится, новый уже будет готов к эксплуатации. Проблема только в том, что у нас выделение памяти операционной системой может быть долгим. И его разбить на отдельные операции уже не получится, так что в таком случае мы не можем сделать по-честному. А ещё такой способ, увы, не универсальный. Мы не всегда знаем точно, когда у нас произойдёт долгая операция. Тут мы уже предвидели, что будет делаться долго, и начали делать это по чуть-чуть на каждом шаге.
    \paragraph{\undercolorblack{orange}{Связный список}.}
    Сначала посмотрим на вычислительную модель, с которой мы будем работать. Она называется <<pointer machine>>. В RAM всё лежало в большом массиве данных. Тут всё лежит в <<узлах>>. Узел --- это класс с какими-то полями. И эти поля могут содержать данные, а могут --- указатели на другие узлы. И полей у каждого узла константное число. За одну операцию можно \begin{enumerate}[1)]
        \item Создать новый узел.
        \item Удалить узел.
        \item Взять значение поля.
        \item Изменить значение поля.
    \end{enumerate}
    Также у нас есть $O(n)$ локальных переменных. Мы можем несложно симулировать pointer machine в RAM-модели. А для того, чтобы сделать наоборот нам не хватает массивов. И зачем нам тогда это ваша pointer machine? За счёт того, что она упрощённая, с ней проще работать. В C/C++ можно работать с памятью руками, а в Java --- нет, специально чтобы было проще. Итак, у нас нет массивов. Зато мы можем сделать список. Для этого мы делаем узлы, в каждом из котором хранятся данные и указатель на следующий узел. Последний узел указывает в \mintinline{java}{null}. А ещё у нас есть начало списка --- первый узел. Мы можем очень легко это перебрать:
    \begin{minted}{\pseudo}
        iterate(func):
            x = head
            while x != null
                func(x)
                x = x.next
    \end{minted}
    Это работает за $O(n)$, какая неожиданность. Что мы не умеем? Мы не умеем быстро достать $k$-тый элемент. Но что мы умеем делать другого? Очень легко вставлять элемент, к тому же в любое место. В массиве делать это очень сложно. А тут --- вообще на изи:
    \begin{minted}{\pseudo}
        add(x, y):
            y.next = x.next
            x.next = y
    \end{minted}
    Эта штука вставляет \pseudocode{y} после \pseudocode{x}. Удалять элементы из списка сложнее. Мы не знаем предыдущий элемент. Есть 2 пути это фиксить. Первый --- удалять не текущий элемент, а следующий. То есть если мы удаляем \pseudocode{x}, то нам нужно передать на вход предыдущий. Второй --- это двусвязный список. Мы просто в каждый узел помимо \pseudocode{next} вставляем \pseudocode{prev}.
    \begin{minted}{\pseudo}
        remove(x):
            if x.next == null && x.prev == null
                head = tail = null
            else if x.prev == null
                head = head.next
                x.prev = null
            else if x.next == null
                tail = tail.next
                x.next = null
            else
                x.next.prev = x.prev
                x.prev.next = x.next
    \end{minted}
    Выглядит как кринж. Это можно исправить очень легко. Давайте добавим в начало и в конец фиктивные элементы. И их мы никак не будем добавлять и удалять, они всегда есть. Тогда останутся только последние 2 строчки. И с добавлением была бы такая же фигня, если бы не фиктивные элементы. Кстати, можно сделать ровно один фиктивный элемент, закольцевав список. Это ничего не усложнит, то даст нам $O(1)$ памяти.\\
    Нам когда-нибудь пригодится структура данных, являющаяся кучкой множеств, с которыми можно делать следующее: добавить элемент в множество, удалить уже найденный элемент и взять произвольный элемент в конкретном множестве. Понятно, что представив множества как наборы двусвязных циклических списков, мы всё получим. Также мы очень легко сможем склеивать два множества (подразумевая, что они не пересекаются). Просто берём фиктивные элементы, через них мы знаем начало и конец всего, перекидываем указатели и чиллим.\\
    Вернёмся к стекам и очередям. Двусвязный список --- это дек. А для стеков и очередей будет достаточно односвязного списка. В стеке \pseudocode{next} --- это указатель от верхних элементов к нижним.
    \begin{minted}{\pseudo}
        push(x):
            x.next = top
            top = x

        pop():
            res = top
            top = top.next
            return res
    \end{minted}
    В очереди указатели \pseudocode{next} указывают в направлении от головы к хвосту.
    \subparagraph{\undercolorblack{orange}{Персистентные структуры данных}.}
    Что хочется? Была СД. Мы делаем с ней что-то, после чего старое состояние умирает. А персистентная СД помнит все свои состояния. Есть разные уровни персистентности. Можно пойти в старую версию и попросить что-то, что не меняет СД. Это частичная персистентность. А есть ещё полная. Там можно откатиться назад и сделать там любое действие. Ни кой это надо, мы поговорим в следующем семестре. А сейчас просто сделаем персистентный стек. Давайте при добавлении (\pseudocode{push}) не только добавим элемент, а ещё и делаем новый указатель на начало. Собственно, это новое начало будет началом следующей <<версии>>. И когда мы откатываемся к предыдущей версии и добавляем в неё, мы снова делаем новое начало, которое имеет следующим элементом начало именно этой версии. И \pseudocode{pop} тоже понятно, как работает. Точно также. Вообще если в pointer machine на каждый узел указывает не более, чем константа узлов, то можно сделать из этой СД персистентную. С очередью нельзя сделать так легко, как со стеком. У неё есть то место, когда добавляются новые элементы. И добавив два элемента в одной и той же версии, мы по-разному изменим указатель у последнего элемента. Можно попытаться сделать очередь на двух стеках, но там возникает проблема: если возвращаться к структуре данных с большим потенциалом, то амортизационное время помрёт. Поэтому мы хотим сделать очередь из двух стеков (понятно, персистентных) без амортизации. Здесь должно быть описание, каким именно образом она делается, но я сам в объяснении Павла Юрьевича нихуя не понимаю, чтобы об этом писать, так что если вам очень нужна эта структура данных, правильнее всего будет посмотреть на \href{https://neerc.ifmo.ru/wiki/index.php?title=\%D0\%9F\%D0\%B5\%D1\%80\%D1\%81\%D0\%B8\%D1\%81\%D1\%82\%D0\%B5\%D0\%BD\%D1\%82\%D0\%BD\%D0\%B0\%D1\%8F_\%D0\%BE\%D1\%87\%D0\%B5\%D1\%80\%D0\%B5\%D0\%B4\%D1\%8C}{Викиконспектах}.
    \paragraph{\undercolorblack{orange}{Система непересекающихся множеств}.}
    Это структура, в которой есть две операции, \pseudocode{union} и \pseudocode{find}. У нас есть какие-то элементы и в каждый момент времени они объединены в какие-то множества. \pseudocode{find(x)} возвращает множество, в котором содержится \pseudocode{x}. \pseudocode{union(s1, s2)} объединяет два множества в одно. Чтобы было проще, делают так: чтобы не создавать дополнительный тип объектов (множество), каждому множеству дают <<главный элемент>> и по ним уже определяют множества. Соотвественно, \pseudocode{union} принимает на вход два таких <<главных элемента>>. Как реализовывать это? Будем хранить каждое множество в виде произвольного подвешенного дерева с корнем в том самом <<главном элементе>>. Как делать \pseudocode{find}? Очень легко, идём вверх, пока не придём в корень. Простоты ради, кстати, можно в конях делать ссылки на самих себя.
    \begin{minted}{\pseudo}
        find(x):
            if parent(x) == x
                return x
            return find(parent(x))
    \end{minted}
    Как делать \pseudocode{union}? Ну, тоже можно очень легко:
    \begin{minted}{\pseudo}
        union(x, y):
            parent(x) = y
    \end{minted}
    Ещё можно в начале сделать \pseudocode{x = find(x)} и \pseudocode{y = find(y)} на случай, если нам не <<главные элементы>> передали. Ещё нужно ничего не делать, если \pseudocode{x = y}.\\
    В чём проблема всего этого? В скорости. Мы можем в \pseudocode{union} построить бамбук, в котором \pseudocode{find} будет работать за $O(n)$. Есть два способа это пофиксить.
    \subparagraph{\undercolorblack{orange}{Ранговая эвристика}.}
    В каждой вершине \pseudocode{x} будем хранить \pseudocode{rank(x)}. Если нам дали деревья с корнями разных рангов, то мы меньшее подвешиваем к большему. Если дали с одинаковыми, то подвешиваем как угодно, но у нового <<главного>> увеличиваем ранг на 1. То есть по сути \pseudocode{rank} --- это самый глубокий путь от листа до корня в дереве.
    \begin{minted}{\pseudo}
        union(x, y):
            if rank(x) > rank(y)
                swap(x, y)
            parent(x) = y
            if rank(x) = rank(y)
                rank(y)++
    \end{minted}
    Утверждение: в такой модификации высота дерева (то есть именно $\rank$) не больше, чем $\log n$. Как мы это докажем? Покажем, что в дереве с корнем в \pseudocode{x} хотя бы $2^{\pseudocode{rank(x)}}$ элементов. Если мы подвесили маленькое дерево к большому, то у нас количество вершин увеличивается, а ранг не меняется, подходит, а если два одинаковых дерева, то в обоих из них $\geqslant2^N$ вершин, значит суммарно $\geqslant 2^{N+1}$, очень хорошо сочетается с увеличенным рангом. Утверждение: во всех деревьях вместе вершин ранга \pseudocode{r} не больше $\frac n{2^{\pseudocode r}}$. Почему? Ну, у нас в каждом поддереве больше либо равно $2^{\pseudocode r}$ вершин, а они не вложенные. А значит, если мы делаем ранговую эвристику, то у нас время работы \pseudocode{find} --- $O(n)$.
    \subparagraph{\undercolorblack{orange}{Эвристика сжатия путей}.}
    Пусть мы запустили \pseudocode{find(x)}. И мы дошли до корня. Если мы запустим \pseudocode{find(x)} ещё раз, мы сделали то же самое. Так зачем это делать ещё раз, если можно сказать, что теперь \pseudocode{x} подвешен к корню. И всем его предкам можно то же самое делать.
    \begin{minted}{\pseudo}
        find(x):
        if parent(x) != x
            parent(x) = find(parent(x))
        return parent(x)
    \end{minted}
    Мы хотим доказать нечто страшное. Если мы делаем сжатие путей, то у нас $O\left(\frac{\log\frac{n^2}m}{\log\frac mn}\right)$. Если мы делаем и то, и другое, то у нас получается $\alpha(n;m)$ --- обратная функция Аккермана. $m$ --- это число операций. У нас после каждого долгого \pseudocode{find} дерево улучшается. Отсюда и формула $\frac{\log\frac{n^2}m}{\log\frac mn}$. Что эта штука говорит? Если $m\approx n$, то там получится $\frac{\log n}{\mathrm{const}}$. А если $m\approx n\sqrt n$, то получится $\frac{\log\sqrt n}{\log\sqrt n}=1$. То есть чем больше запросов, тем меньше формула. Верно, что $O\left(\frac{\log\frac{n^2}m}{\log\frac mn}\right)=O(\log n)$, но $O(\log n)$ --- это скучно и не даёт много информации. Итак, давайте докажем страшную формулу. Разобьём все рёбра на классы. Класс ребра из вершины \pseudocode{x} в вершину \pseudocode{parent(x)} --- это максимальное \pseudocode{i}, что $\pseudocode{rank(parent(x))}-\pseudocode{rank(x)}\geqslant b^{\pseudocode{i}}$, где $b$ --- константа. Правильно брать её $\frac mn$. Но классы у нас будут не от 0 до $\infty$, а до $Z$, где $Z$ --- это и есть $\frac{\log\frac{n^2}m}{\log\frac mn}$. То есть классы больше $Z$ мы приравниваем к $Z$. Кстати, все формулы тут ломаются на маленьких числах, но мы можем просто везде насыпать прибавление небольших констант, чтобы всё работало. Мы так уже сделали, когда смотрели на асимптотику при $m\approx n$. Но мы забьём на это. Посчитаем время на все операции. Возьмём рёбра, по которым мы прошли в \pseudocode{find}'е. Посмотрим на их классы. Ну, они какие-то. Хочется понять, сколько рёбер мы суммарно прошли за несколько \pseudocode{find}'ов. Давайте в каждом классе выберем <<самое верхнее ребро>> на пути каждого запроса? Сколько таких рёбер может быть? Ну, явно не больше $Z$. То есть у нас $mZ$ таких рёбер всего. Давайте посчитаем <<не-верхние рёбра>> для всех классов от 0 до $Z-1$ (то есть кроме последнего). Возьмём ребро (\pseudocode{x}\textrightarrow\pseudocode{parent(x)}). Оно было в классе $k$. Значит разница рангов между \pseudocode{x} и \pseudocode{parent(x)} больше либо равно $b^k$. И где-то потом выше в том же запросе мы тоже встретили ребро класса $k$. Сколько времени ребро может быть не-последним класса $k$? На самом деле не очень долго. Заметим, что после сжатия пути разница рангов между \pseudocode{x} и \pseudocode{parent(x)} (новым \pseudocode{parent}) --- как минимум $2b^k$, потому что у нас хотя бы два раза ранг увеличился на не меньше, чем $b^k$ (один раз на \pseudocode{x}\textrightarrow\pseudocode{parent(x)} и ещё один раз на верхнем ребре класса $k$). Ранги мы, кстати, не трогаем при сжатии. Они при сжатии сохраняются. То есть мы к разнице рангов приделываем $b^k$. То есть после $b$ раз класс ребра совершенно точно увеличивается. То есть у нас есть $n$ рёбер, по каждому мы идём $b$ раз, после чего оно переходит в следующий класс. Соответственно, чтобы перевести его в класс $Z$, нужно $nbZ$ проходов. А теперь утверждается, что рёбер класса $Z$ бывает не очень много. Каждое из них означает разность рангов $\geqslant b^Z$. У нас ранги не больше $n$, значит на каждом пути встречается не больше $\frac n{b^Z}$ таких рёбер. Значит всего их не больше $\frac{n^2}{b^Z}$. То есть суммарно мы прошли $mZ+nbZ+\frac{n^2}{b^Z}$. Взяв $b=\frac mn$ и $Z=\frac{\log\frac{n^2}m}{\log\frac mn}$, получим, что эта сумма --- это $mZ+nnZ+m=O(mZ)$. То есть амортизированное время --- $O(Z)$, то есть то, что нужно.
    \subparagraph{\undercolorblack{orange}{Функция Аккермана}.}
    Построим таблицу. В начальной строке находятся числа 0, 2, 4, 6, 8,... В начальном столбце --- нули, в следующем --- двойки, а все остальные элементы считаются так: \pseudocode{A[i][j] = A[i-1][A[i][j-1]]}. Нулевая строка --- это $2x$. Первая --- это $2^x$ (мы в индексе \pseudocode{i} мы берём \pseudocode{i}-тый элемент предыдущей строки). Во второй строке мы (после 0 и 2) берём $2^2$, $2^{2^2}$, $2^{2^{2^2}}$,... Короче, что-то очень быстрорастущее. В следующей строке вообще получается сначала 4, потом $2^{16}$ (так как это четвёртое число предыдущей строки), а потом что-то большое и страшное. Нужен элемент предыдущей строки с номером $2^{16}$.\\
    Хорошо, а что такое обратная функция Аккермана? А это следующее: $\alpha(m;n)=\min k:\pseudocode{A[k][m/n]}\geqslant\log n$ (опять эта формула не работает для маленьких значений, но нам всё ещё пофиг). Теперь докажем, что применение обоих эвристик даёт нам именно $O(\alpha(m,n))$. Опять вводим класс ребра. Теперь это $\min k:\exists i~\pseudocode{A[k][i]}\leqslant\pseudocode{rank(x)}<\pseudocode{rank(parent(x))}\leqslant\pseudocode{A[k][i+1]}$. Опять ограничим это число сверху числом $Z=\alpha(m,n)$. Опять посчитаем, сколько рёбер мы встретили. <<Последних рёбер>>, как обычно, $mZ$. Теперь <<не-последние>> на классах от 0 до $Z-1$. Пусть есть 2 ребра из класса $k$ и при этом ребро мы рассматриваем не-последнее из них. Оно имеет класс $k$, то есть какое-то число Аккермана \pseudocode{A[k-1][j1]} находится между рангами \pseudocode{rank(x)} и \pseudocode{rank(parent(x))}. С <<последим ребром>> класса $k$ та же фигня, там есть число \pseudocode{A[k-1][j2]}. То есть после каждой операции сжатия пути между рангами \pseudocode{rank(x)} и \pseudocode{rank(parent(x))} появляется новое \pseudocode{A[k-1][j2]} из предыдущей строчки функции Аккермана. Сколько их нужно, чтобы и в текущей строчке ранги не поместились? \pseudocode{A[k][j]} хватит, где \pseudocode{A[k][j]} --- то значение функции Аккермана, что \pseudocode{rank(x)} и \pseudocode{rank(parent(x))} помешались между \pseudocode{A[k][j]} и \pseudocode{A[k][j+1]}. Вершин с рангом на отрезке между \pseudocode{A[k][j]} и \pseudocode{A[k][j+1]} может быть $\sum_{r=\pseudocode{A[k][j]}}^{\pseudocode{A[k][j+1]}}\frac n{2^r}$ Каждое $\frac n{2^r}$ не больше $\frac{2n}{2^{\pseudocode{A[k][j]}}}$. Для каждой из таковых мы сделали $\pseudocode{A[k][j]}$ операций, а $\sum_j\pseudocode{A[k][j]}\frac{2n}{2^{\pseudocode{A[k][j]}}}\leqslant2n\sum_i\frac i{2^i}=n\cdot\mathrm{const}$. То есть всего (на $Z$ классов мы опать получили $Zn$ операций).\\
    На не-последних рёбрах класса $Z$ мы делаем $\frac mn$ операций на каждого, то есть $m$ всего. Если $\frac mn$ немножко растёт, то у нас получается асимптотика такая, как нам нужно.
    \paragraph{\undercolorblack{orange}{Фиббоначева куча}.}
    Итак, в куче мы умели удалять минимум и вставлять элемент. Мы всё это делали за $O(\log n)$. А ещё мы искали минимум за $O(1)$. Мы хотим ещё пару операций:
    \begin{itemize}
        \item Объединение двух куч (\pseudocode{merge}).
        \item Уменьшение ключа (\pseudocode{decrease_key}). Так получается, что алгоритмы монотонны, поэтому нужно только уменьшать значения, а не увеличивать их.
    \end{itemize}
    Фиббоначева куча делает
    \begin{itemize}
        \item \pseudocode{insert} за единицу.
        \item \pseudocode{get_min} за единицу.
        \item \pseudocode{merge} за единицу.
        \item \pseudocode{decrease_key} за единицу.
    \end{itemize}
    Если бы она ещё и минимуи за единицу удаляла, мы бы сортировали массив за $O(n)$.\\
    Начнём с биномиальной кучи. Из неё сделаем фиббоначеву.
    \paragraph{\undercolorblack{orange}{Биномиальные куча}.}
    Двоичные деревья --- это слишком просто, вместо него биномиальные. Дерево B\textsubscript0 --- одна вершина:
    \begin{center}
        \begin{tikzpicture}[
            scale=.5
            ]
            \node[treenode] (0) at (0,0) {};
        \end{tikzpicture}
    \end{center}
    B\textsubscript1 --- это две вершины, родитель и потомок:
    \begin{center}
        \begin{tikzpicture}[
            scale=.5
            ]
            \node[treenode] (0) at (0,0) {};
            \node[treenode] (1) at (0,-2) {};
            
            \draw[->] (0) -- (1);
        \end{tikzpicture}
    \end{center}
    Вот так выглядит B\textsubscript2:
    \begin{center}
        \begin{tikzpicture}[
            scale=.5
            ]
            \node[treenode] (0) at (0,0) {};
            \node[treenode] (1) at (0,-2) {};
            \node[treenode] (2) at (2,-2) {};
            \node[treenode] (21) at (2,-4) {};
            
            \draw[->] (0) -- (1);
            \draw[->] (0) -- (2);
            \draw[->] (2) -- (21);
        \end{tikzpicture}
    \end{center}
    А так --- B\textsubscript3:
    \begin{center}
        \begin{tikzpicture}[
            scale=.5
            ]
            \node[treenode] (0) at (0,0) {};
            \node[treenode] (1) at (0,-2) {};
            \node[treenode] (2) at (2,-2) {};
            \node[treenode] (3) at (4,-2) {};
            \node[treenode] (21) at (2,-4) {};
            \node[treenode] (31) at (4,-4) {};
            \node[treenode] (32) at (6,-4) {};
            \node[treenode] (321) at (6,-6) {};
            
            \draw[->] (0) -- (1);
            \draw[->] (0) -- (2);
            \draw[->] (0) -- (3);
            \draw[->] (2) -- (21);
            \draw[->] (3) -- (31);
            \draw[->] (3) -- (32);
            \draw[->] (32) -- (321);
        \end{tikzpicture}
    \end{center}
    Каждое следующее (B\textsubscript{$k$}) строится так: мы берём предыдущее дерево (B\textsubscript{$k-1$}) и к его корню подвешиваем ещё одно такое же. Несложно понять, что в дереве $2^k$ вершин. У корня всегда $k$ вершин.\\
    Как на этих деревьях делать кучу? По стандартному правилу, предок меньше потомка. А если у нас не степень двойки элементов? То берём несколько биномиальных деревьев. То есть 6 элементов --- это два дерева: B\textsubscript1 и B\textsubscript2. Как считать минимум? Пробежаться по всем $\log n$ деревьям и искать минимум в корнях. Можно сохранить ссылку на него, если хочется не $O(\log n)$, а $O(1)$. Как выполнять все действия?
    \begin{itemize}
        \item[\pseudocode{merge}] Это просто сложение битовых чисел. Если у мы склеиваем два дерева одного размера, мы склеиваем их, а результат --- \textit{переносим через разряд}. Как мы склеиваем два дерева одинакового размера? Очень просто, к тому, у которого корень меньше, подвесить второе. Это делается за $O(1)$, а значит весь \pseudocode{merge} --- за $O(\log n)$. Понятно, как пересчитывать минимум. Задача. Как хранить дерево, чтобы к нему было легко что-то подвешивать? Да легко, храним детей каждой вершины в виде списка.
        \item[\pseudocode{insert}]
        Ну, сливаем данную кучу с B\textsubscript0. Это будет даже амортизированно за $O(1)$, как с двоичным счётчиком. 
        \item[\pseudocode{remove_min}]
        Все дети минимума образуют биномиальную кучу. Берём её и сливаем со всем остальным.
        \item[\pseudocode{decrease_key}]
        Берём и просеиваем ключ вверх. Это работает за $O(\log n)$.
    \end{itemize}
    Вот иллюстрация к \pseudocode{merge}. Пусть у нас есть два дерева на 6 и 3 вершины соответственно:
    \begin{center}
        \begin{tikzpicture}[
            scale=.5,framed
            ]
            \node[treenode] (0) at (-3,0) {3};
            \node[treenode] (1) at (-3,-2) {7};
            \draw[->] (0) -- (1);
            
            \node[treenode] (0) at (0,0) {1};
            \node[treenode] (1) at (0,-2) {6};
            \node[treenode] (2) at (2,-2) {4};
            \node[treenode] (21) at (2,-4) {8};
            
            \draw[->] (0) -- (1);
            \draw[->] (0) -- (2);
            \draw[->] (2) -- (21);
        \end{tikzpicture}
        \hspace{5mm}
        \begin{tikzpicture}[
            scale=.5,framed
            ]
            \node[treenode] (0) at (-3,0) {13};
            
            \node[treenode] (0) at (0,0) {2};
            \node[treenode] (1) at (0,-2) {10};
            \draw[->] (0) -- (1);
        \end{tikzpicture}
    \end{center}
    Тогда маленькое деревце (13) мы оставляем as is, а два дерева B\textsubscript1 нужно склеить. Получится B\textsubscript2, которое у нас тоже уже есть, в связи с чем их тоже нужно склеить. Получится такое:
    \begin{center}
        \begin{tikzpicture}[
            scale=.5,framed
            ]
            \node[treenode] (A0) at (-3,0) {13};
            
            \node[treenode] (0) at (0,0) {1};
            \node[treenode] (1) at (0,-2) {6};
            \node[treenode] (2) at (2,-2) {4};
            \node[treenode] (3) at (4,-2) {8};
            \node[treenode] (21) at (2,-4) {2};
            \node[treenode] (31) at (4,-4) {10};
            \node[treenode] (32) at (6,-4) {3};
            \node[treenode] (321) at (6,-6) {7};
            
            \draw[->] (0) -- (1);
            \draw[->] (0) -- (2);
            \draw[->] (0) -- (3);
            \draw[->] (2) -- (21);
            \draw[->] (3) -- (31);
            \draw[->] (3) -- (32);
            \draw[->] (32) -- (321);
        \end{tikzpicture}
    \end{center}
    Вернёмся к фиббоначивой куче. В чём проблема с биномиальной? Все деревья слишком одинаковы, нет пространства для манёвра. Фиббоначева куча устроена похоже на биномиальную, но там деревья страшные.\\
    Вершина ранга $k$ выглядит так: у неё есть $k$ детей. При этом самый левый ребёнок имеет класс хотя бы 0, следующий --- хотя бы 1, последний --- хотя бы $k-1$. В биномиальной было \textbf{ровно} 0, 1, ..., $k-1$, а тут \textbf{хотя бы}. Ещё мы разрешим укороченную вершину ранга $k$ (обозначаем её $k^*$), которая выглядит идентично вершине ранга $k-1$, но она ведёт себя чуть-чуть по-другому. Мы можем подвесить $k^*$ туда, где нужно хотя бы $k$, а $k-1$ --- не можем. Корни, понятно, укороченными быть не могут. Например, дерево может выглядеть так:
    \begin{center}
        \begin{tikzpicture}[scale=.5]
            \node[treenode] (0) at (0,0) {\bf3};
            \node[treenode] (1) at (0,-2) {\bf1};
            \node[treenode] (2) at (2,-2) {\bf1*};
            \node[treenode] (3) at (4,-2) {\bf2*};
            \node[treenode] (11) at (0,-4) {\bf0};
            \node[treenode] (31) at (4,-4) {\bf2};
            \node[treenode] (311) at (4,-6) {\bf0};
            \node[treenode] (312) at (6,-6) {\bf1*};
            
            \draw[->] (0) -- (1);
            \draw[->] (0) -- (2);
            \draw[->] (0) -- (3);
            \draw[->] (1) -- (11);
            \draw[->] (3) -- (31);
            \draw[->] (31) -- (311);
            \draw[->] (31) -- (312);
        \end{tikzpicture}
    \end{center}
    Самое главное про это дерево --- ранг дерева равен $O(\log n)$. Докажем это. Достаточно доказать, что зависимость количества вершин от ранга растёт экспоненциально. Посмотрим на самое маленькое дерево ранга $k$. Это, понятно, $k^*$, у которого все дети минимальные и со звёздочкой: 0, $1^*$, $2^*$, ..., $(k-1)^*$. Несложно по индукции доказать, что в такого рода деревьях количество вершин --- это $k$-тое число Фиббоначи. А $F_n=O\left(\left(\frac{1+\sqrt5}2\right)^n\right)$.\\
    И вот фиббоначева куча --- это набор таких деревьев и ссылка на минимум. Этот кринж очень легко \pseudocode{merge}'ить. У нас было 2 набора деревьев, мы сделаем из них один. И минимум пересчитаем. Как делать \pseudocode{insert}? Через \pseudocode{merge}. А \pseudocode{get_min}? Смотрим по ссылке. Остался \pseudocode{remove_min}. Мы берём детей данного минимума, объединяем со всеми остальными деревьями и чиллим... Нет! Нам нужно минимум найти. А хрен мы его найдём. Поэтому вводится операция \pseudocode{compress}, которая ищет минимум и заодно сжимает одинаковые вершины. Просто идём слева направо. Имеем массив, в котором по индексу храним, где мы встречали дерево с рангом таким индексом. Когда мы находим второе дерево с таким же рангом, мы можем их как биномиальные деревья приклеить друг к другу. Понятно, что размер массива --- $O(\log n)$, потому что ранги $O(\log n)$. Понятно, что после такой операции мы можем найти минимум за $\log n$, у нас именно столько деревьев. Утверждается, что сжатие работает амортизировано быстро. Кладём монетку на корень каждого дерева. Когда мы сливаем два дерева (не кучи), у нас из двух деревьев получается одно, а монетку мы тратим именно на слияние.\\
    А как же делать \pseudocode{decrease_key}? Просто просеивать нельзя, ведь бамбук --- это правильное дерево. Что же мы делаем, если мы уменьшили элемент, и он стал меньше родителя? Мы напрочь отрезаем его вместо со всеми потомками и просто запихиваем его в наш список деревьев. Если у родителя не было звёздочки, мы ставим её и побеждаем. А что, если эта звёздочка уже была? Берём этого родителя, говорим, что у него ранг $k-2$ (без звёздочки) и помещаем в список деревьев. Делаем так для всех родителей, пока не доходим до корня (ему просто уменьшаем ранг) или пока не найдём кого-нибудь без звезды. А в каком месте это $O(1)$? Тут мы вводим другой потенциал --- число деревьев + удвоенное количество звёздочек. На все операции, кроме \pseudocode{decrease_key} обновление потенциала не влияет, а тут мы проходим по $m$ вершинам, и как раз уменьшаем потенциал на $m$ (и суммарно увеличиваем на константу, но не суть).\\
    Это всё в теории очень быстро, но константа там очень большая, так что на практике не используется.
    \paragraph{\undercolorblack{orange}{Левосторонняя куча}.}
    Давайте хранить кучу в немножко скошенном двоичном дереве. Для каждой вершины храним число --- максимально близкое свободное место снизу. Утверждение --- это число не больше логарифма количества вершин в поддереве. Это очевидно. Как добавлять элемент? Идём до ближайшего пустого места (вправо) и туда подвешиваем. Правда, свойство нарушается. Просто берём и меняем детей местами, если надо. После чего просеиваем всё вверх.
    Как склеивать кучи? Ну, кто корень? Минимальный из корней. Берём правое поддерево меньшего и рекурсивно склеиваем. После этого меняем местами детей, если надо. Понятно, это $\log(n_1+n_2)$. Удаление минимума делается через склеивание детей корня.
    \pseudocode{decrease_key} делается так: мы отрезаем нужную вершину, а потом обратно сливаем. Правда, перед слитием нужно пересчитать числа в вершинах всех предков удалённой вершины, но это легко, идём вверх и выбираем меньшее из номера в этом нашем хождении и того, что уже было написано в вершине.
    \paragraph{\undercolorblack{orange}{Динамическое программирование}.}
    Давайте посчитаем числа Фибоначчи. $F_n=F_{n-1}+F_{n-2}$. Напишем процедуру, которая из считает.
    \begin{minted}{\pseudo}
        fib(n):
            if n <= 2
                return 1
            else
                return fib(n - 1) + fib(n - 2)
    \end{minted}
    Это работает дохрена долго (за $O(F_n)$ и работает), потому что в \pseudocode{fib(n - 1)} мы вычисляем \pseudocode{f(n - 2)}, то есть его мы считаем дважды. А меньшие штуки ещё большее количество раз считаются. Выглядит как биба. Идея в том, чтобы каждое число считать один раз. Пусть 1:
    \begin{minted}{\pseudo}
        F = [-1] * (n + 1)
        fib(n):
            if F[n] != -1
                return F[n]
            if n <= 2
                F[n] = 1
            else
                F[n] = fib(n - 1) + fib(n - 2)
            return F[n]
    \end{minted}
    Теперь каждое число Фибоначчи считается не больше одного раза, что, безусловно, неплохо, то есть этот код работает за $O(n)$.
    Но на самом деле можно делать всё то же самое, но проще. Мы точно знаем, в каком порядке мы считаем значения. От меньшего к большему. Значит всё хорошо считается циклом.
    \begin{minted}{\pseudo}
        F[1] = F[2] = 1
        for i in [3, n]
            F[i] = F[i - 1] + F[i - 2]
        return F[n]
    \end{minted}
    На самом деле тут можно ещё и сэкономить память. Если мы вычисляем число Фибоначчи только 1 раз, то можно заменить это всё так:
    \begin{minted}{\pseudo}
        a = b = 1
        for i in [3, n]
            a, b = b, a + b
        return b
    \end{minted}
    Если можно использовать это только для странных функций, то скучно. Поэтому поехали в тупые примеры.\\
    Есть полоска из клеточек. На самой левой стоит кузнечик. Ему нужно допрыгать до самой правой клетки. При этом он умеет прыгать либо на 1, либо на 2 клетки. Хочется узнать, сколькими способами можно допрыгать с первой клетки до последней. Как это считать? Да легко. Как мы могли придти в некую клетку? Либо из предыдущей, либо из пред-предыдущей. Причём оба этих варианта не пересекаются. То есть количество путей всего --- это сумма для этих двух предыдущих клеток. То есть мы опять получаем $d(n)=d(n-1)+d(n-2)$, где $d(n)$ --- количество способов дойти до клетки с номером $n$. Несложно понять, что $d(0)=d(1)=1$. То есть это опять Фибоначчи. Что изменится, если кузнечик может и на 3 клеточки прыгать? Тогда $d(n)=d(n-1)+d(n-2)+d(n-3)$. А если на $k$? То же самое, но тут задавать первые $k$ условий сложно. Поэтому тут пишется так:
    \begin{minted}{\pseudo}
        d[0] = 1
        for i in [1, n)
            d[i] = 0
            for j in [1, k]
                if i - j >= 1
                    d[i] += d[i - j]
        return d[n]
    \end{minted}
    За сколько работает? За $O(nk)$. Как оптимизировать? Можно хранить сумму предыдущих $k$ чисел, чтобы её не пересчитывать. И тогда внутренний цикл уничтожится, а на его месте появится
    \begin{minted}{\pseudo}
        d[i] = s
        s += d[i]
        if i - k >= 1
            s -= d[i - k]
    \end{minted}
    Обычно такие задачи довольно искусственно выглядит. Менее искусственные динамические задачи --- оптимизация. Например, если кузнечик что-то собирает или что-то теряет.\\
    У Георгия Александровича новая игра. Вы можете делать домашки через одну (то есть две соседние сделанные должны быть либо последовательными, либо между ними должна быть ровно одна ещё). За каждую известно, сколько баллов за оформление кода теряет студент, если её сделает. Цель: сделать домашки так, чтобы потерять минимальное количество баллов. Как это делать? Да точно также. $d(n)$ --- это минимальная сумма на пути до домашки номер $n$. Либо вы делали предыдущую домашку, либо пред-предыдущую, и так далее. Тут получается $d(n)=\min{d(n-1);d(n-2)}$.
    \begin{minted}{\pseudo}
        d[1] = c[1]
        d[2] = c[2]
        for i in [3, n]
            d[i] = min(d[i - 1], d[i - 2]) + cnt[i]
        return F[n]
    \end{minted}
    Проблема в том, что непонятно, а какие именно домашки решать. Наша программа скажет нам, сколько минимум баллов мы потеряем, но не скажет, как. Чтобы это поправить, можно в каждом хранить предыдущее, из которого мы должны были придти, чтобы получить оптимальное данное.
    \begin{minted}{\pseudo}
        for i in [3, n]
            if d[i - 1] < d[i - 2]
                d[i] = d[i - 1]
                prev[i] = i - 1
            else
                d[i] = d[i - 2]
                prev[i] = i - 2
    \end{minted}
    Но это ест память. А можно не хранить, а просто идти с конца и просто смотреть, отличаются ли $d(n)$ и $d(n-i)$ на $cnt(n)$? Если да, значит мы пришли в $d(n)$ именно из $d(i)$.
    \begin{minted}{\pseudo}
        result = {n}
        i = 0
        while i != 1
            if d[i] = cnt[i] + d[i - 1]
                result.add(i - 1)
                i = i - 1
            else
                result.add(i - 2)
                i = i - 2
        return result
    \end{minted}
    Но тут мы по сути пишем одну и ту же программу дважды, что неудобно, если вычисление функции долгое.
    А ещё можно делать динамику назад. У нас есть наше текущее состояние $n$ и мы перебираем, все места, куда мы можем из него попасть и берём более оптимальное из того, что там было и того, что мы хотим записать. Например, только так и делается, если мы из клетки $i$ можем прыгнуть в клетку $i+j$, где $j$ --- делитель $i$.\\
    По сути, то же самое можно делать на многомерном пространстве. Например, на клетчатой доске. Если мы идём по ней только вправо или вниз, то более или менее понятно, мы для каждой клетки считаем, сколько способов попасть клетку сверху и в клетку слева. И берём минимум/сумму/небо/Аллаха, в зависимости от того, что в задаче хотят. Надо только считать клеточки в том порядке, чтобы предыдущие точно были посчитаны. И точно также можно делаyouть динамику назад, почему нет.
    \paragraph{\undercolorblack{orange}{Фибоначчи за логарифм}.}
    Как считать числа Фибоначчи на $O(\log n)$? Поскольку растут эти числа быстро (экспоненциально), один хрен ты посчитаешь их честно, они выводится в ответ будут за $O(n)$. Поэтому лучше считать их по какому-нибудь модулю. Мы храним два числа, $F_i$ и $F_{i+1}$. Мы должны получить следующую пару. Это линейное преобразование, а значит можно выразить его умножением на матрицу. Спойлер: эта матрица будет $\left(\begin{matrix}
        0&1\\
        1&1
    \end{matrix}\right)$.
    То есть $\left(\begin{matrix}
        F_{n-1}\\
        F_n
    \end{matrix}\right)=\left(\begin{matrix}
        0&1\\
        1&1
    \end{matrix}\right)^{n-2}\left(\begin{matrix}
        1\\
        1
    \end{matrix}\right)$
    А матрицы можно возводить в степень за при помощи \pseudocode{fast_pow} за $O(\log n)$.
    \paragraph{\undercolorblack{orange}{Расстояние Левинштейна}.}
    Поехали в более жизненные задачи про динамику. Есть такая утилита, \mintinline{bash}{diff}. Она сравнивает два файлика, и говорит, какие строчки появились, какие --- удалились. Работает она при помощи ДП. Решим более удобную задачу. У нас есть две строки, и хотим понять минимальное количество изменений, чтобы из первой получить вторую. Например, <<\textbf{полено}>> и <<\textbf{племя}>>. <<Изменение>> --- это либо добавить букву, либо убрать, либо заменить. Тут мы убираем обе буквы <<\textbf{о}>>, а <<\textbf{н}>> заменяем на <<\textbf{м}>>. То есть ответ --- 3. Строим динамику так: \pseudocode{d[i;j]} --- это количество действий, чтобы из \pseudocode{A[0:i)} получить \pseudocode{B[0:j)}. Как это считать? Если кто-то из \pseudocode{i} и \pseudocode{j} равен нулю, то нам нужно просто добавить/убрать все буквы, то есть ответ --- тот, кто нулю не равен. Если \pseudocode{A[i-1]} и \pseudocode{B[j-1]} --- одно и то же, то \pseudocode{d[i;j]} равно \pseudocode{d[i-1;j-1]}. Если не равны, нужно как-то изменить конец префикса. Можно заменить букву, и будет ответ \pseudocode{d[i-1;j-1]+1}. Можно убрать букву в первом слове (\pseudocode{A}) и пытаться получить тот же префикс, но из более маленького префикса \pseudocode{A}. Это \pseudocode{d[i-1;j]+1}. Либо же можно добавить букву во второе слово, получив \pseudocode{d[i;j-1]+1}. Надо взять из этих трёх значений минимум. Посчитав это, получим такую таблицу
    \begin{center}
        \begin{tabular}{c|ccccccc}
            \hline
            &&П&О&Л&Е&Н&О\\
            \hline
            &0&1&2&3&4&5&6\\
            П&1&0&1&2&3&4&5\\
            Л&2&1&1&1&2&3&4\\
            Е&3&2&2&2&1&2&3\\
            М&4&3&3&3&2&2&3\\
            Я&5&4&4&4&3&3&3\\
            \hline
        \end{tabular}
    \end{center}
    Ответ --- 3. Как получить сами действия? Как всегда.\\
    За сколько это работает? За произведение длин. Это много. Что делает \mintinline{bash}{diff}, чтобы работать быстрее? Большая часть таблицы не нужна, потому что ему обычно скармливают похожие файлы. Что можно делать? Можно хранить только какие-то маленькие значения этой таблицы. Обычно (для похожих файлов) маленькие числа стоят в районе диагонали.
    \paragraph{\undercolorblack{orange}{Текстовый редактор}.}
    Есть текстовые редакторы, которым хочется разбить ваш длинный текст по строкам. Если переносить жадно, в конце строк могут образовываться пробелы. Это некрасиво. Как этот вопрос решён в \TeX? У нас есть слова (каждое с шириной \pseudocode{a[i]}), и есть ширина страницы \pseudocode{W}. Что значит, что в конце нет больших пробелов? Некрасивость пробела в конце равна кубу его длины. Мы минимизируем сумму этих кубов. При этом нельзя переносить слова. \pseudocode{d[i]} --- минимальный штраф, чтобы расставить первые \pseudocode{i} слов. Тогда перебираем такие \pseudocode{k}, что последние \pseudocode{k} слов можно разместить на последней строке. Соответственно, берём ответ для \pseudocode{d[i-k]} и складываем с кубом пробела в конце строки (если последние \pseudocode{k} строк помещаются в \pseudocode{W}). И потом берём минимум.
    \paragraph{\undercolorblack{orange}{Перемножение матриц}.}
    Если у нас есть $A\in\scriptK^{l\times n}$ и $B\in\scriptK^{n\times m}$. Их перемножение --- $C\in\scriptK^{l\times m}$. Как совсем втупую их перемножать? За $O(mnk)$. А как перемножить $k$ штук матриц? Ну, по ассоциативности можно как угодно. Мы и хотим решить, в каком порядке. Пусть \pseudocode{d[l,r]} --- это время, за которое перемножается отрезок матриц \pseudocode{A[l;r)}. Если в массиве один элемент, то мы считаем произведение за 0 действий. Если нет, давайте посмотрим, где у нас последнее произведение. Ну, где-то на полуинтервале. При этом перемножаем мы уже 2 готовых отрезка. То есть мы перебираем \pseudocode{m} и нам нужно минимизировать сумму \pseudocode{d[l,m]}, \pseudocode{d[m;r]} и произведения размеров \pseudocode{l}-той, \pseudocode{m}-той и \pseudocode{r}-той матрицы. Проблема в том, в каком порядке считать отрезки. Самый простой вариант --- отрезки по возрастанию длины. То есть мы перебираем сначала длину, а потом левую границу, и норм. А можно перебирать левую границу в обратном порядке, а правую --- в прямом. Тогда к тому моменту, как мы дойдём до \pseudocode{d[l;r]}, мы точно подсчитаем все подотрезки.
    \paragraph{\undercolorblack{orange}{Ускорение трёхмерных динамик}.}
    Мы хотим разбить массив на \pseudocode{m} кусков и минимизировать квадрат суммы всех кусков. Считаем такую динамику: \pseudocode{d[i;j]} --- разбиение первых \pseudocode{i} элементов на \pseudocode{j} кусков. Тогда \pseudocode{d[i;j]} --- \pseudocode{min(d[k;j-1]+S(k;i)^2)}. При этом \pseudocode{d[0;j]} --- \pseudocode{0}. Ну, дальше понятно.
    \begin{minted}{\pseudo}
        for i in [1:n]
            for j in [1:m]
                for k in [0:i)
                    d[i;j] = min(d[i;j], d[k;j-1]+S(k;i)*S(k;i))
    \end{minted}
    Оптимизация этого: Смотрим на последний цикл. Пы перебираем, где идеально резать последний кусок. Давайте запомним это в \pseudocode{P[i;j]} (то есть запишем туда \pseudocode{k}). И что? Посмотрим на последний элемент. Утверждается, что при разбиении всего без него на столько же кусков, граница не будет левее. То есть что \pseudocode{P[i;j]>=P[i-1][j]}. Теперь другая ситуация. Пусть мы разбили те же элементы, но на один фрагмент больше. Теперь утверждается, что \pseudocode{P[i;j] <= P[i;j+1]}. Доказывать мы это не будем. А значит мы сильно сжимаем границы перебора \pseudocode{k}. Проблема в том, что мы не посчитали \pseudocode{P[i;j+1]}. Ну, давайте перебирать \pseudocode{j} в обратном порядке. Почему стало лучше? Сколько все эти циклы теперь работают? Нужно посчитать сумму $\sum_{\pseudocode{i},\pseudocode{j}}\pseudocode{P[i;j+1]-P[i-1;j]}$. Говорят, что почти всё сократится. Не сократится только то, что \pseudocode{P[n;j+1]} или \pseudocode{P[i;k+1]}. Их \pseudocode{n} штук. А их сумма --- \pseudocode{n^2}.
    \paragraph{\undercolorblack{orange}{Задача о рюкзаке}.}
    Дан рюкзак, в который можно положить предметы с весом суммарно не больше $S$. Предметы все обладают весом $w_i$ и стоимостью/ценностью $c_i$. Требуется напихать в рюкзак предметы так, чтобы их суммарная ценность была максимально.\\
    Что про эту задачу нужно знать? То, что она NP-полная. Что это, мы изучим потом, для нас сейчас это значит, что один хрен мы решим её за полиномиальное время от количества предметов. А если научимся, у нас полетит криптография. И даже при $w_i=c_i$ задача всё равно NP-полная.\\
    Один из примеров, когда эту задачу можно быстро решить --- маленькое $S$ и целые $w_i$. Тогда задача решается динамикой. Идём по предметам. Каждый предмет мы либо берём, либо нет. Нам нужно помнить, сколько предметов мы уже положили и их суммарный вес. Тогда \pseudocode{d[i][j]} --- можно ли из первых $i$ предметов выбрать набор с суммарным весом $j$. Тогда, если мы берём предмет, то у нас \pseudocode{d[i][j] = d[i - 1][j]}. Если не берём --- \pseudocode{d[i][j] = d[i - 1][j - w[i]]}. То есть если у нас в массиве --- <<возможно ли>>, то мы просто берём логическое <<ИЛИ>> этих двух значений и всё готово. Ответ --- максимальное $j$, у которого \pseudocode{d[n][j]} истинно. Можно, чтобы не писать этот цикл, под \pseudocode{d[i][j]} иметь ввиду вес не ровно $j$, а не больше $j$. Давайте теперь хранить не <<можно ли>>, а максимум стоимости первых $i$ предметов с суммой $j$. Тогда всё меняется не сильно: \pseudocode{d[i][j] = max(d[i - 1][j], d[i - 1][j - w[i]]  + c[i])};\\
    Чтобы хранить такую большую табличку, нужно $O(nS)$ памяти. А мем в том, что используем мы только предыдущую строку, значит можно хранить только $O(S)$ памяти, имеет предыдущую строку, имеем текущую, пересчитываем. А ещё можно хранить вообще только одну строку, а не две. Как? Да легко. Когда мы вычисляем элемент строки, мы используем только всё, что слева от него. Так давайте пересчитывать строку от больших значений к маленьким. Тут проблема с восстановлением пути, мы могли всё перезатереть к концу. Поэтому есть специальная техника. Мы знаем, что наш путь идёт всегда только направо и вниз. В конце мы знаем только последнюю строку. Чтобы найти оптимальный ответ --- ищем максимум в последней строке. Мы хотим узнать, где был наш путь в строке $\frac n2$. Для каждой клетки храним это самый индекс. В строке с номером $\frac n2$ храним тупо номер столбца, а в каждой следующей смотрим, откуда мы пришли и записываем в ответ ту самую клетку $\frac n2$-ной строки. дальше нам нужно найти то, что было до строки $\frac n2$ и то, что после. В процессе динамики мы сохраним среднюю строку. И от начала до неё делаем абсолютно ту же самую динамику. Так мы находим, как мы пришли из строки $\frac n4$ в строку $\frac n2$. И из строки $\frac{3n}4$ в строку $n$ тоже пересчитаем. Каждую из них считаем рекурсивно. То есть мы в динамике сохраняем среднюю и нижнюю строки. Потом перезапускаем динамику 2 раза, от верхней левой четверти и правой нижней. Утверждается, что это работает за линию. Понятно, что всё работает за линейную память. Заметим, что на каждом следующем уровне рекурсии мы отрезаем половину матрицы. Проблема с памятью. Мы от меньшей части матрицы запускаем рекурсию (у нас $\log n$) уровней рекурсии, и в каждом из них массив, который мы храним, вдвое меньше. То есть рекурсия занимает $O(n)$ памяти. А вторую половину запускаем как хвостовую рекурсию.\\
    Другой способ решать рюкзак мы узнаем чуть позже, а сейчас
    \subparagraph{\undercolorblack{orange}{Приближённая задача о рюкзаке}.}
    К вам приходят и говорят, что хочется не обязательно самое оптимальное, а отличающееся от него на $1-\varepsilon$. Такие штуки называются приближёнными алгоритмами, и для многих NP-полных задач такие придуманы. И в некоторых задачах понятно, как решив задачу для некоторого $\varepsilon_0$, то понятно, как довести до оптимального, а в некоторых (как тут) --- вообще нет.\\
    Самый тупой алгоритм --- отсортируем по убыванию предметы по отношению стоимости к весу. Жадно запихаем <<самые ценные>> в рюкзак. Утверждается, что это почти не плохое решение. Когда же происходит всё плохо? Например, когда в рюкзак размером 1000 вы пихаете предмет с ценностью 2 и весом 1 вместо предмета с ценностью 1001 и весом 1000. Вариант 2 --- берём максимум из жадного алгоритма и самого ценного предмета. Утверждается, что для $\varepsilon=0.5$ мы решили задачу. У нас были какие-то предметы. И есть последний не поместившийся. Если посчитать ценность вместе с этим предметом, получится что-то точно большее, чем оптимальный ответ. А значит мы либо берём первые предметы, либо что-то не меньшее, чем последний предмет.\\
    Ещё вариант. Давайте сделаем, чтобы у наших предметов $c_i$ были маленькими и целыми. Скажем, что новые ценности предметов --- это какие-то $c'_i=\lceil\frac{c_i}B\rceil$, где $B$ --- некоторая выбранная нами константа. Как решать такой рюкзак? Динамикой <<для каждой стоимости минимальный вес>>. Среди первых $i$ предметов, разумеется. Это считается за $O(n\sum c'_i)\leqslant O(n^2\lceil\frac{c_{\max}}B\rceil)$. Утверждается, что мы не сильно ошибёмся после преобразования ценностей. Мы посчитали каждую стоимость в попугаях, равных $B$. То есть у каждой стоимости потерялся остаток от деления на $B$. Таких остатков --- $n$, максимальный из размер --- $B$. То есть $nB\leqslant\varepsilon\cdot\mathrm{OPT}$ (OPT --- оптимальное решение). Понятно, что OPT не меньше, чем $c_{\max}$. То есть возьмём $B=\frac{\varepsilon c_{\max}}n$. И всё будет хорошо. Если поставить $B$ в асимптотику, получится $O\left(\frac{n^3}\varepsilon\right)$.
    \paragraph{\undercolorblack{orange}{Динамика по множествам}.}
    Как мы уже выяснили, есть задачи, которые хорошо за полином не решаются. И люди меряются тем, а за что же можно их решать. Начнём с того, что решим рюкзак за $2^n$ на какой-то полином ($n^k$). Это легко, перебираем все варианты <<взять очередной предмет или нет>>. Это $2^n$ вариантов. Перебираем все, смотрим результат, обновляем ответ, если текущий вариант оптимальнее. Как это число технически сделать? Будем кодировать всё $n$-битными числами, где $i$-тый бит равен единице, если мы берём соответствующий предмет. В чём прелесть? В том, что с числами в битовой записи легко работать компьютеру.\\
    Поскольку каждое число --- это множество, нам хочется понять, как выполнять операции на множествах. Как сделать множества из элемента $x\in[0:n-1]$? Это \mintinline{c++}{1 << x}. Объединение множеств $a$ и $b$ --- это \mintinline{cpp}{a | b}, а пересечение --- \mintinline{cpp}{a & b}. А разность --- это чуть сложнее: \mintinline{cpp}{a & (~b)} или \mintinline{cpp}{a ^ (a & b)}. Как проверить, лежит ли элемент $x$ в множестве $a$? Ну, так: \mintinline{cpp}{a & (1 << x)} или так: \mintinline{cpp}{(a >> x) & 1}.\\
    Что можно сделать лучше в нашем переборе? Попилим наш массив предметов пополам, переберём все множества слева и все справа. Отдельно друг от друга переберём. Для каждого первого куска найдём подходящий ему второй кусок. То есть нам нужно подобрать к весу первого куска $w_x$ подобрать второй кусок с максимальным весом $w_y$, удовлетворяющим $w_x+w_y\leqslant S$. Это делается путём сортировки каждого куска (за $\frac n22^{\frac n2}$) и построения префиксных максимумов. Итого было у нас $n2^n$, а стало $n2^{\frac n2}$, что уже очень значимо, константа в показателе экспоненты --- это приятно.\\
    Давайте рассмотрим другую задачу. У вас есть не один рюкзак, а бесконечное количество пакетов, но с меньшей грузоподъёмностью. И предметы нужно все разложить. И пакетов хочется минимум. И тут уже совсем всё плохо, даже если грузоподъёмность пакета $S$ мала. В рюкзаке состояние --- это число, а тут --- не число, а дичь какая-то. Как это решать? \pseudocode{d[x]} --- это минимальные число пакетов, чтобы разложить предметы множества \pseudocode{x}. То есть параметр динамики --- это множество, а не чиселко. Как такое пересчитывать? Есть множество, которое мы раскладываем по пакетам. Нам нужно взять какой-то пакет и свести задачу к меньшей. Переберём то, что мы положим в очередной пакет. То есть мы перебираем \pseudocode{y} --- подмножество \pseudocode{x} такие, которые кладутся в пакет. То есть для каждого такого подмножества ответик --- \pseudocode{1 + d[x - y]}. И надо взять минимум из всех таких. Как это написать? Аккуратно, нам нужно, чтобы к текущему состоянию предыдущие были посчитаны. Ну, можно числа по возрастанию перебирать, подмножество всегда не больше множества.
    \begin{minted}{\pseudo}
        d[0] = 0
        for x in [1, 2 ** n - 1]
            for y in [1, 2 ** (n - 1)]
                if x & y == y and w[y] <= S
                  d[x] = min(d[x], 1 + d[x - y])
    \end{minted}
    Это работает за долго, за $4^n$. Хочется перебирать подмножества проще, не беря все множества в принципе и не проверяя на подмножественность. Как их так перебирать? По убыванию. Самое первое --- само \pseudocode{x}. Следующее --- самую младшую единичку установить в 0. Потом --- следующую в ноль, а предыдущую --- в 1. То есть в каждый момент мы убираем последнюю единичку, а все предыдущие возвращаем в 1. Это же можно сделать так: вычитаем единицу из подмножества, а потом делаем побитовое <<И>> с множеством. Первый пункт уберёт нудный нам ноль, всё ниже задаст единичку, но нам из них нужны только те, что были в \pseudocode{x}.
    \begin{minted}{\pseudo}
        d[0] = 0
        for x in [1, 2 ** n - 1]
            y = x
            while y > 0
                if x & y == y and w[y] <= S
                    d[x] = min(d[x], 1 + d[x - y])
                y = (y - 1) & x
    \end{minted}
    Это работает за $3^n$. Хорошо, но хочется ещё ускорить. Почему работает долго? Потому что переход динамики сложный, мы кладём в рюкзак сразу много предметов. Вместо этого давайте оставим последний пакет незавязанным. И хранить мы будем не одно число (количество пакетов), а ещё и вес последнего незаполненного пакета. И минимум мы будем брать так: меньше пакетов --- хорошо, а если их одинаковое количество, то хорошо --- когда последний пакет легче. Почему? Действительно, почему же если пакетов больше, но последний пакет легче --- это ситуация лучше, чем если пакетов меньше, но вес последнего --- больше? Ну, во втором мы можем просто завязать последний пакет, получить количество пакетов либо равное, либо всё ещё меньшее первого случая, а последний пакет пуст. Как писать динамику на этом?
    \begin{minted}{\pseudo}
        for x in [1, 2 ** n - 1]
            for i in [0, n]
                if x & (1 << i)
                    a, b = d[x - 1 << i]
                    if b + w[i] <= S
                        a, b = a, b + w[i]
                    else
                        a, b = a + 1, w[i]
                    d[x] = min(d[i], {a, b})
    \end{minted}
    \paragraph{\undercolorblack{orange}{Динамика по профилю}.}
    Задача: есть прямоугольничек $n\times m$. Мы хотим забить его доминошками (если кто-то не знает, что это, во-первых, позор ему, во-вторых, прямоугольник $1\times2$). Требуется посчитать количество способов сделать это. Небезызвестный товарищ Кохась научился делать это, считая определки упоротых матриц, но мы сейчас будем делать не это, а вообще частный случай задачи, когда $m$ маленькое. Тогда будем идти по столбцам и считать, что у нас из $i$-того столбца вылезают в правую половину. Нам нужно помнить, какие. Это мы опять же кодируем битовым числом (длины $m$). И у нас по сути появляется ломаная, которую и называют профилем в названии <<Динамика по профилю>>. То есть наша динамика --- это \pseudocode{d[i][p]} --- число способов заполнить первые \pseudocode{i} столбцов, чтобы вылез профиль \pseudocode{p}. Как это пересчитывать? Нам дан профиль, и мы хотим заполнить новый столбец. У нас есть пустые клетки. Если дырка в профиле имеет размер 1, то выбора у нас нет, мы ставим туда горизонтальную доминошку.  Если дырка имеет размер 2, мы либо ставим туда одну доминошку вертикально, либо две горизонтально. Но вообще перебирать для всех возможных размеров дырок, это сложно.
    \begin{minted}{\pseudo}
        for i in [0, n - 1]
            for p in [0, 2 ** m - 1]
                for q --- подходящий к р новый профиль
                    d[i + 1][q] += d[i][p]
    \end{minted}
    Два цикла (по \pseudocode{p} и по \pseudocode{q}) немножко бесят. Проблема опять в том, что мы добавляем сразу много доминошек. Наш переход добавляет целый слой, а мы не хотим целый слой. У нас можно добавлять по одной доминошке за раз. У нас будет заполнена какая-то левая часть таблицы, но на не ровная, а в ней есть уголочек, первые $j$ клеток в ней заполнены. И там ещё профиль. То есть у нас заполнены первые $i$ столбцов целиком, из самого $i$-го столбца заполнены первые $j$ клеток, а потом из этого всего выпирает профиль $p$. То есть динамика теперь трёхмерная. Если в профиле выставлен $j$-й бит, значит наша клеточка $(i,j)$ и так уже заполнена, можно ничего не делать. Если там 0, мы либо кладём доминошку вертикально (не не всегда так можем), либо горизонтально. Когда мы заполним последнюю клетку столбца, мы должны скопировать результаты текущего столбца с $j=m$ в результат следующего с $j=0$.
    \begin{minted}{\pseudo}
        for i in [0, n - 1]
            for j in [0, m - 1]
                for p in [0, 2 ** m - 1]
                    if p & (1 << j)
                        d[i][j + 1][p - (1 << j)] += d[i][j][p]
                    else
                        d[i][j + 1][p + (1 << j)] += d[i][j][p]
                        if p & (1 << (j + 1)) == 0
                            d[i][j + 2][p] += d[i][j][p]
            for p in [0, 2 ** m - 1]
                d[i + 1][0][p] = d[i][m][p]
    \end{minted}
    Ответ будет в \pseudocode{d[n][0][0]}.
    \paragraph{\undercolorblack{orange}{Хэш-таблицы}.}
    Очень популярная СД, все ими пользуются даже не зная об этом. Зачем нам хэш-таблицы? В основном для структур типа map и структур типа set. Кто такой set? Это штука, в которую можно добавлять элементы (\pseudocode{insert(x)}), удалять элементы (\pseudocode{remove(x)}) и проверять, есть ли заданный элемент в ней (\pseudocode{contains(x) -> bool}). Кто такой map? По жизни --- соответствие из каких-то ключей в какие-нибудь значения. И в ней есть \pseudocode{put(k, v)} --- добавить по ключу значение и \pseudocode{get(k) -> v} --- взять значение по ключу (если есть, иначе \pseudocode{null}). Тривиально можно сделать map из set.\\
    Как сделать map тривиально (если ключи маленькие и притом целые)? На массиве, \pseudocode{put} добавляет по индексу значение, а \pseudocode{get} берёт его.\\
    А если ключи большие и целые, то как жить? Идея в следующем: заводить большой массив плохо, поэтому делаем функцию, которая из большого числа делает маленькое, (в заданном интервале). Она и называется хэш-функцией. Например, можно взять $f\colon x\mapsto x\operatorname{mod} M$. Тут есть проблема, если у нас ключи с одинаковым хэшем, мы при \pseudocode{get} возвращаем не то, если мы делаем \pseudocode{get} не по тому значению которое записали, но с тем же хэшем (например, при $N=4$ комбинация \pseudocode{put(1, v)}, \pseudocode{get(5)} будет фигня, надо вернуть \pseudocode{null}, а мы вернём значение и даже не поймём, что что-то не так). И это исправляется легко, в массиве просто храним пару из самого ключа и значения. Но другая проблема: мы перезаписываем значения с разными ключами, но одинаковыми хэшами. Это называтся <<коллизия>>, и так легко уже не фиксится. Но давайте посмотрим, а можно ли сделать вероятность коллизии маленькой? Давайте брать абсолютно случайную хэш-функцию. Пусть у нас функция возвращает значения от 0 до M. Посчитаем математическое ожидание количества коллизий. Каждая пара образует коллизию с вероятностью $\frac1M$. Значит мат. ожидание --- $\frac{n(n-1)}{2M}\approx\frac{n^2}M$. Если хочется, чтобы коллизий было мало это ожидание должно быть порядка 1, а значит размер массива нужно брать $n^2$. Это хорошо, но как-то очень далеко.\\
    Поэтому с коллизиями смиряются. Как жить с коллизиями? Например, в массиве хранить списочек всех ключей и значений. Тогда так выглядит код:
    \begin{minted}{\pseudo}
        put(k, v):
            a[h(k)].push_back({k, v})
        get(k):
            for k1, v in a[h(k)]
                if k1 == k
                    return v
             return null
    \end{minted}
    Тут в \pseudocode{put} есть проблема, если ключи одинаковые добавлять, но это сами как-нибудь пофиксите. За сколько этот код работает? Если у нас плохая хэш-функция, то это линейное время. Если мы считаем, что функция хорошая и вообще всё равномерно распределяет, то это работает за $O(\frac nm)$. Тогда мы делаем $m\approx n$ и всё хорошо, лучше не получится. Но такую функцию мы делать не можем, увы. Заметим, что придумать ровно одну хорошую функцию мы не можем, ведь у неё будут фиксированные коллизии. Поэтому давайте брать случайную из некоторого семейства.\\
    Например, это можно сделать так. Мы генерируем случайное большое простое число $P$ и случайное число $A$ от 1 до $P-1$. Тогда утверждается, что семейство \pseudocode{f(k): (k * A) \% P \% M
    }.
    Чему равно мат. ожидание? По всем элементам просуммируем вероятность того, что какой-то другой попадёт в ту же ячейку. Утверждается, что она порядка $\frac1m$. Что значит, что у нас совпали хэши? Это значит, что $k_1A\operatorname{mod}P\operatorname{mod}M=k_2A\operatorname{mod}P\operatorname{mod}M$. Это можно записать как $(k_1-k_2)A\operatorname{mod}P\operatorname{mod}M=0$. Вероятность того, что $k_1$ и $k_2$ равны по модулю $P$ крайне мала в силу того, что $P$ --- большое и простое. В противном случае нам нужно, чтобы $(k_1-k_2)A$ попало в какое-то из чисел $0$, $M$, $2M$, ... --- $\frac PM$. Для каждого из них есть только одно возможное $A$, потому что по простому модулю можно не только умножать на $A$, но и делить. Поскольку $A$ случайное, вероятность выбрать плохое $A$ --- $\frac1M$. А как раз $\frac1M$ мы и получали. А значит мы снова получаем, что мат. ожидание количества совпавших элементов --- $O(\frac nm)$.\\
    Для разных структур данных нужны разные свойства хэшей, для хэш-таблицы --- только вероятность коллизии в $\frac1M$. Иногда кажется, что нам хватит какого-то другого свойства. Например, достаточно ли нам, что вероятность того, что $h(x)=y$ равна $\frac1M$. Нет, не достаточно, мы можем взять семейство функций <<константа от $1$ до $M$>>. Тогда взяв случайную, мы и получим вероятность $\frac1M$. Но все элементы попадают в одно место, это точно очень плохо.\\
    Финальный пункт: у нас время $O(1)$ при линейной памяти. Но у нас массив списков, с этим очень неудобно жить. Поэтому есть второй метод --- открытая адресация. Заведём массив на побольше (на $2N$ элементов, а не $N$, про запас). У нас некоторые клеточки будут заняты, а некоторые --- свободны. И теперь вместо того, чтобы хранить список, при добавлении элемента идти <<вправо>> до нахождения ближайшей свободной клетки (дошли до конца --- переходим в начало). Как искать? Да легко, идём пока не встретим либо подходящую, либо пустую. В первом случае возвращаем её, иначе --- \pseudocode{null}. Есть одна большая проблема --- удаление элементов. Понятно, в чём проблема, будут дырки, и искать так, как мы искали, не получится. Понятен инвариант, для любого ключа всё от его ячейки до того, где он находится на самом деле, заполнено. Есть тупой вариант, элементы мы не удаляем до конца, а забиваем их мусором, когда мусора много --- перехэшируем (см. ниже). Но это не очень хорошо. Давайте нормально. Мы удаляем элемент. Это может кому-то другому сломать инвариант. Давайте от места удаления идти направо и смотреть, нет ли нарушения. Если оно есть, мы перемещаем в дырку этот самый элемент. Но он образовывает новую дырку. Так пойдём снова вправо с новой дыркой. Когда мы дошли до пустого места, перестаём идти. Вся эта штука (открытая адресация в принципе) очень хорошо кэшируется, потому что мы идём по массиву подряд. Оценивать асимптотику этого долго и сложно. Это мы делать не будем. Вроде доказано, что размер $2N$ достаточен, но мы не будем об этом говорить.\\
    Хорошо, давайте о следующем поговорим. Мы сделали всё по ТЗ, выбрали функцию, понапихали элементы, а потом осознали, что нас предали. Поэтому можно выкинуть хэш-функцию, взять новую и за линейное время сделать новую хэш-таблицу из тех же элементов. Обычно это помогает. Но есть проблема. А ещё обычно хэш-таблицы не знают, сколько в неё кладут. Ну так давайте делать резиновый массив, который мы также, как и в \mintinline{cpp}{std::vector} увеличивать в два раза.
\end{document}